\subsection{Linear Algebra}

Linear algebra is the branch of mathematics that studies vectors, matrices, and linear transformations. It provides the foundation for solving systems of equations and analysing multidimensional data. Its applications span across machine learning, computer graphics, physics, finance, and engineering.\\
References from \cite{garcia2017second} and \cite{axler2024linear}.

\subsubsection{Vector Spaces}

\begin{definition} \hlt{Vector Space}\\
Let $\mathbb{F} = \mathbb{R}$ or $\mathbb{C}$. Then $\mathcal{V}$ is a vector space over a field $\mathbb{F}$ if the following are satisfied:
\begin{enumerate}[label=\roman*.]
\setlength{\itemsep}{0pt}
\item Identities: $\exists \ \mathbf{0} \in \mathcal{V}$ such that $\mathbf{0} + \mathbf{u} = \mathbf{u} \ \forall \ \mathbf{u} \in \mathcal{V}$, and $\exists \ \mathbf{1} \in \mathcal{V}$ such that $\mathbf{1} \mathbf{u} = \mathbf{u} \ \forall \ \mathbf{u} \in \mathcal{V}$
\item Commutativity: $\mathbf{u} + \mathbf{v} = \mathbf{v} + \mathbf{u} \ \forall \ \mathbf{u}, \mathbf{v} \in \mathcal{V}$
\item Associativity: $\mathbf{u} + (\mathbf{v} + \mathbf{w}) = (\mathbf{u} + \mathbf{v}) + \mathbf{w}$ and $(ab) \mathbf{u} = a(b \mathbf{u}) \ \forall \ a, b \in \mathbb{F}$ and $\forall \ \mathbf{u}, \mathbf{v}, \mathbf{w} \in \mathcal{V}$
\item Distributive: $a(\mathbf{u} + \mathbf{v}) = a \mathbf{u} + a \mathbf{v}$ and $(a + b)\mathbf{u} = a \mathbf{u} + b \mathbf{u} \ \forall \ a, b \in \mathbb{F}$ and $\forall \ \mathbf{u}, \mathbf{v} \in \mathcal{V}$
\item Additive Inverse: $\forall \ \mathbf{u} \in \mathcal{V} \ \exists$ unique vector $-\mathbf{u} \in \mathcal{V}$ such that $\mathbf{u} + (- \mathbf{u}) = \mathbf{0}$
\item Multiplicative Inverse: $\forall \ \mathbf{u} \in \mathcal{V}$ with $\mathbf{u} \neq \mathbf{0} \ \exists$ unique $\mathbf{w} \in \mathcal{V}$ such that $\mathbf{u} \mathbf{w} = 1$
\end{enumerate}
\end{definition}

\begin{remark} \hlt{Properties of Vector Spaces}\\
Let $\mathcal{V}$ be a vector space over a field $\mathbb{F}$.
\begin{enumerate}[label=\roman*.]
\setlength{\itemsep}{0pt}
\item A vector space has a unique additive identity
\item Every element in a vector space has a unique additive inverse
\item $\mathbf{0} \mathbf{v} = \mathbf{v}$ for every $\mathbf{v} \in \mathcal{V}$
\item $a\mathbf{0} = \mathbf{0}$ for every $a \in \mathbb{F}$
\item $(-1)\mathbf{v} = -\mathbf{v}$ for every $\mathbf{v} \in \mathcal{V}$
\end{enumerate}
\end{remark}

\begin{definition} \hlt{Subspace}\\
A subspace of an $\mathbb{F}$-vector space $\mathcal{V}$ is a subset $\mathcal{U}$ of $\mathcal{V}$ that is an $\mathbb{F}$-vector space with the same vector addition and scalar multiplication operations as in $\mathcal{V}$. To check $\mathcal{U}$ satisfies the following:
\begin{enumerate}[label=\roman*.]
\setlength{\itemsep}{0pt}
\item Additive Identity: $\mathbf{0} \in \mathcal{U}$
\item Closure Under Addition: $\mathbf{u}, \mathbf{v} \in \mathcal{U}$ implies $\mathbf{u} + \mathbf{v} \in \mathcal{U}$
\item Closed Under Scalar Multiplication: $a \in \mathbb{F}$ and $\mathbf{u} \in \mathcal{U}$ implies $a \mathbf{u} \in \mathcal{U}$.
\end{enumerate}
\end{definition}

\begin{definition} \hlt{Sum of Subspaces}\\
Suppose $\mathcal{U}_1, \ldots, \mathcal{U}_m$ are subspaces of $\mathcal{V}$. The sum of $\mathcal{U}_1, \ldots, \mathcal{U}_m$ as denoted by $\mathcal{U}_1 + \cdots + \mathcal{U}_m$ is defined to be set of all possible sums of elements of $\mathcal{U}_1, \ldots, \mathcal{U}_m$:
\begin{align}
\mathcal{U}_1 + \cdots + \mathcal{U}_m \coloneqq \{ \mathbf{u}_1 + \cdots + \mathbf{u}_m : \mathbf{u}_1 \in \mathcal{U}_1, \ldots, \mathbf{u}_m \in \mathcal{U}_m \} \nonumber
\end{align}
\end{definition}

\begin{definition} \hlt{Direct Sum of Subspaces}\\
If $\mathcal{U}_1, \ldots, \mathcal{U}_m$ are subspaces of $\mathcal{V}$ such that $\mathcal{V} = \mathcal{U}_1 + \cdots + \mathcal{U}_m$, then every element of $\mathcal{V}$ may be written as
\begin{align}
\mathbf{u}_1 + \cdots + \mathbf{u}_m \nonumber
\end{align}
where each $\mathbf{u}_j \in \mathcal{U}_j$. Then $\mathcal{V}$ is the direct sum of subspaces $\mathcal{U}_1, \ldots, \mathcal{U}_m$, written as $\mathcal{V} = \mathcal{U}_1 \oplus \cdots \oplus \mathcal{U}_m$
\end{definition}

\begin{definition} \hlt{Linear Combination}\\
Let $\mathcal{U}$ be non-empty subset of $\mathbb{F}$-vector space $\mathcal{V}$. A linear combination of elements of $\mathcal{U}$ is expressed as
\begin{align}
c_1 \mathbf{v}_1 + c_2 \mathbf{v}_2 + \cdots + c_r \mathbf{v}_r \nonumber
\end{align}
where $r$ is a positive integer, and $\mathbf{v}_1, \mathbf{v}_2, \ldots, \mathbf{v}_r \in \mathcal{U}$, $c_1, c_2, \ldots, c_r \in \mathbb{F}$.\\
A linear combination is \hlt{trivial} if $c_1 = c_2 = \cdots = c_r = 0$.
\end{definition}

\begin{definition} \hlt{Span}\\
Let $\mathcal{U}$ be subset of $\mathbb{F}$-vector space $\mathcal{V}$. If $\mathcal{U} \neq \emptyset$, then $\text{span}(\mathcal{U})$ is the set of linear combinations of elements of $\mathcal{U}$.\\
Define $\text{span}(\emptyset) = \{\mathbf{0}\}$. If $\beta = \mathbf{v}_1, \mathbf{v}_2, \ldots, \mathbf{v}_r$ is a list of $r \geq 1$ vectors in $\mathcal{V}$, define $\text{span}(\beta) = \text{span}\{ \mathbf{v}_1, \mathbf{v}_2, \ldots, \mathbf{v}_r \}$.
\end{definition}

\begin{lemma} \label{lem:lineardependencelemma} \hlt{Linear Dependence Lemma}\\
If $(\mathbf{v}_1, \ldots, \mathbf{v}_m)$ is linearly dependent in $\mathcal{V}$ and $\mathbf{v}_1 \neq 0$, then $\exists \ j \in \{2, \ldots, m \}$ such that the following hold:
\begin{enumerate}[label=\roman*.]
\setlength{\itemsep}{0pt}
\item $\mathbf{v}_j \in \text{span}(\mathbf{v}_1, \ldots, \mathbf{v}_{j-1})$
\item if the $j$th term is removed from $(\mathbf{v}_1, \ldots, \mathbf{v}_m)$, the span of the remaining list equals $\text{span}(\mathbf{v}_1, \ldots, \mathbf{v}_m)$.
\end{enumerate}
\begin{proof}
\begin{enumerate}[label=\roman*.]
\setlength{\itemsep}{0pt}
\item Suppose $(\mathbf{v}_1, \ldots, \mathbf{v}_m)$ is linearly dependent in $\mathcal{V}$, and $\mathbf{v}_1 \neq 0$. Then $\exists \ a_1, \ldots, a_m \in \mathbb{F}$, not all $0$, s.t.
\begin{align}
a_1 \mathbf{v}_1 + \cdots + a_m \mathbf{v}_m = 0 \nonumber
\end{align}
Not all of $a_2, \ldots, a_m$ can be $0$ as $\mathbf{v}_1 \neq 0$. Let $j$ be largest element of ${2, \ldots, m}$ s.t. $a_j \neq 0$. Then
\begin{align}
\mathbf{v}_j = - \frac{a_1}{a_j} \mathbf{v}_1 - \cdots - \frac{a_{j-1}}{a_j} \mathbf{v}_{j-1} \nonumber
\end{align}
\item Suppose $\bm{u} \in \text{span}(\mathbf{v}_1, \ldots, \mathbf{v}_m)$. Then $\exists \ c_1, \ldots, c_m \in \mathbb{F}$ such that
\begin{align}
\mathbf{u} = c_1 \mathbf{v}_1 + \cdots + c_m \mathbf{v}_m \nonumber
\end{align}
Replace $\mathbf{v}_j$ with equation from (i), hence (ii) holds.
\end{enumerate}
\end{proof}
\end{lemma}

\begin{theorem} In a finite-dimensional vector space, length of every linearly independent list of vectors is less than or equal to length of every spanning list of vectors.
\begin{proof} Suppose $(\mathbf{u}_1, \ldots, \mathbf{u}_m)$ is linearly independent in $\mathcal{V}$, and $(\mathbf{w}_1, \ldots, \mathbf{w}_n)$ spans $\mathcal{V}$. Need to prove $m \leq n$.
\begin{enumerate}[label=\roman*.]
\setlength{\itemsep}{0pt}
\item List $(\mathbf{w}_1, \ldots, \mathbf{w}_n)$ spans $\mathcal{V}$, hence adjoining any vector to it produces linearly dependent list. By Lemma \ref{lem:lineardependencelemma}, can remove one of $\mathbf{w}$ such that the list containing $\mathbf{u}_1$ and remaining $\mathbf{w}$'s spans $\mathcal{V}$.
\item List of length $n$ from previous step $j-1$ spans $\mathcal{V}$, hence adjoining any vector to it produces a linearly dependent list. The list of length $(n+1)$ obtained by adjoining $\mathbf{u}_j$ to the list, placing it just after $\mathbf{u}_1, \ldots, \mathbf{u}_{j-1}$ is linearly dependent. By Lemma \ref{lem:lineardependencelemma}, one of the vectors in this list is span of previous ones, because $(\mathbf{u}_1, \ldots, \mathbf{u}_j)$ is linearly independent, hence this vector must be one of $\mathbf{w}$ not one of $\mathbf{u}$. Hence can remove $\mathbf{w}$ from the list so that the new list consists of $\mathbf{u}_1, \ldots, \mathbf{u}_j$ and remaining $\mathbf{w}$ spans $\mathcal{V}$.
\item At step $m$, all $\mathbf{u}$ has been added, and process tops. If at any step $\mathbf{u}$ is added and no more $\mathbf{w}$ could be removed, there is a contradiction. hence there must be at least as many $\mathbf{w}$ as $\mathbf{u}$.
\end{enumerate}
\end{proof}
\end{theorem}

\begin{proposition} Every subspace of finite-dimensional vector space is finite dimensional.
\begin{proof}
Suppose $\mathcal{V}$ is finite dimensional, and $\mathcal{U}$ is subspace of $\mathcal{V}$.
\begin{enumerate}[label=\roman*.]
\setlength{\itemsep}{0pt}
\item If $\mathcal{U} = \{0\}$, then $\mathcal{U}$ is finite dimensional. If $\mathcal{U} \neq \{0\}$, choose nonzero vector $\mathbf{v}_1 \in \mathcal{U}$.
\item If $\mathcal{U} = \text{span}(\mathbf{v}_1, \ldots, \mathbf{v}_{j-1})$, then $U$ is finite dimensional. If $\mathcal{U} \neq \text{span}(\mathbf{v}_1, \ldots, \mathbf{v}_{j-1})$, choose vector $\mathbf{v}_j \in \mathcal{U}$ such that $\mathbf{v}_j \neq \text{span}(\mathbf{v}_1, \ldots, \mathbf{v}_{j-1})$
\end{enumerate}
After each step, a linearly independent list is constructed by Lemma \ref{lem:lineardependencelemma}. This linearly independent list cannot be longer than any spanning list of $\mathcal{V}$, hence process must terminate, which means $\mathcal{U}$ is finite dimensional.
\end{proof}
\end{proposition}

\begin{definition} A \hlt{basis} of $\mathcal{V}$ is a list of vectors in $\mathcal{V}$ that is linearly independent and spans $\mathcal{V}$.
\end{definition}

\begin{proposition} A list $(\mathbf{v}_1, \ldots, \mathbf{v}_n)$ of vectors in $\mathcal{V}$ is a basis of $\mathcal{V}$ if and only if every $\mathbf{v} \in \mathcal{V}$ can be written uniquely in form $\mathbf{v} = a_1 \mathbf{v}_1 + \cdots + a_n \mathbf{v}_n $, where $a_1, \ldots, a_n \in \mathbb{F}$.
\begin{proof}
First suppose $(\mathbf{v}_1, \ldots, \mathbf{v}_n)$ is basis of $\mathcal{V}$. Let $\mathbf{v} \in \mathcal{V}$. As $(\mathbf{v}_1, \ldots, \mathbf{v}_n)$ spans $\mathcal{V}$, $\exists \ a_1, \ldots, a_n \in \mathbb{F}$ such that proposition holds. To show the representation is unique, suppose $b_1, \ldots, b_n$ are scalars, hence
\begin{align}
\mathbf{v} = b_1 \mathbf{v}_1 + \cdots + b_n \mathbf{v}_n \nonumber
\end{align}
Subtracting this from proposition equation,
\begin{align}
0 = (a_1 - b_1) \mathbf{v}_1 + \cdots + (a_n - b_n) \mathbf{v}_n \nonumber
\end{align}
This implies each $a_j - b_j = 0$ as $(\mathbf{v}_1, \ldots, \mathbf{v}_n)$ is linearly independent, hence $a_1 = b_1, \ldots, a_n = b_n$.\\
For other direction, suppose every $\mathbf{v} \in \mathcal{V}$ can be written uniquely in proposition form. To show $(\mathbf{v}_1, \ldots \mathbf{v}_n)$ is linearly independent, suppose $a_1, \ldots a_n \in \mathbb{F}$ are such that
\begin{align}
0 = a_1 \mathbf{v}_1 + \cdots + a_n \mathbf{v}_n \nonumber
\end{align}
The uniqueness of representation from proposition with $\mathbf{v} = 0$ implies $a_1 = \cdots = a_n = 0$. Hence $(\mathbf{v}_1, \ldots, \mathbf{v}_n)$ is linearly independent and hence is a basis of $\mathcal{V}$.
\end{proof}
\end{proposition}

\begin{theorem} \label{thm:spanlistreduce} Every spanning list in a vector space can be reduced to a basis of vector space.
\begin{proof} Suppose $(\mathbf{v}_1, \ldots, \mathbf{v}_n)$ spans $\mathcal{V}$. To remove some vectors from $(\mathbf{v}_1, \ldots \mathbf{v}_n)$ so that remaining vectors form a basis of $\mathcal{V}$. Start with $\mathcal{B} = (\mathbf{v}_1, \ldots, \mathbf{v}_n)$.
\begin{enumerate}[label=\roman*.]
\setlength{\itemsep}{0pt}
\item If $\mathbf{v}_1 = 0$, delete $\mathbf{v}_1$ from $B$, If $\mathbf{v}_1 \neq 0$, leave $\mathcal{B}$ unchanged.
\item If $\mathbf{v}_j$ is in $\text{span}(\mathbf{v}_1, \ldots, \mathbf{v}_{j-1})$, delete $\mathbf{v}_j$ from $\mathcal{B}$. If $\mathbf{v}_j$ is not in $\text{span}(\mathbf{v}_1, \ldots, \mathbf{v}_{j-1})$, leave $\mathcal{B}$ unchanged.
\end{enumerate}
Stop process after step $n$, getting a list $\mathcal{B}$ that spans $\mathcal{V}$, and discarded only vectors that were already in span of previous vectors. Hence $\mathcal{B}$ is linearly independent by Lemma \ref{lem:lineardependencelemma}. Hence $\mathcal{B}$ is a basis of $\mathcal{V}$.
\end{proof}
\end{theorem}

\begin{corollary} Every finite-dimensional vector space has a basis.
\begin{proof} By definition, a finite-dimensional vector space has a spanning list. By Theorem \ref{thm:spanlistreduce}, any spanning list can be reduced to a basis. 
\end{proof}
\end{corollary}

\begin{theorem} Every linearly independent list of vectors in a finite-dimensional vector space can be extended to a basis of the vector space.
\begin{proof} Suppose $\mathcal{V}$ is finite dimensional and $(\mathbf{v}_1, \ldots, \mathbf{v}_m)$ is linearly independent in $\mathcal{V}$. To extend $(\mathbf{v}_1, \ldots, \mathbf{v}_m)$ to a basis of $\mathcal{V}$. Let $(w_1, \ldots, w_n)$ be any list of vectors in $\mathcal{V}$ that spans $\mathcal{V}$.
\begin{enumerate}[label=\roman*.]
\setlength{\itemsep}{0pt}
\item If $\mathbf{w}_1 \in \text{span}(\mathbf{v}_1, \ldots, \mathbf{v}_m)$, let $\mathcal{B} = (\mathbf{v}_1, \ldots, \mathbf{v}_m)$. If $\mathbf{w}_1$ is not in $\text{span}(\mathbf{v}_1, \ldots, \mathbf{v}_m)$, let $\mathcal{B} = (\mathbf{v}_1, \ldots, \mathbf{v}_m, \mathbf{w}_1)$.
\item If $\mathbf{w}_j \in \text{span}(B)$, leave $\mathcal{B}$ unchanged. If $\mathbf{w}_j$ not in $\text{span}(B)$, extend $\mathcal{B}$ by adjoining $\mathbf{w}_j$ to it.
\end{enumerate}
After each step, $\mathcal{B}$ is still linearly independent as otherwise Lemma \ref{lem:lineardependencelemma} would give a contradiction. After step $n$, the span of $\mathcal{B}$ includes all the $\mathbf{w}$. Hence the $\mathcal{B}$ obtained after step $n$ spans $\mathcal{V}$ and hence is a basis of $\mathcal{V}$.
\end{proof}
\end{theorem}

This theorem can be applied to show that every subspace of a finite-dimensional vector space can be paired with another subspace to form a direct sum of the whole space.

\begin{proposition} Suppose $\mathcal{V}$ is finite dimensional and $\mathcal{U}$ is a subspace of $\mathcal{V}$. Then there is a subspace $\mathcal{W}$ of $\mathcal{V}$ such that $\mathcal{V} = \mathcal{U} \oplus \mathcal{W}$.
\begin{proof} As $\mathcal{V}$ is finite dimensional, so is $\mathcal{U}$, hence $\exists \ \text{basis}(\mathbf{u}_1, \ldots, \mathbf{u}_m)$ of $\mathcal{U}$. As $(\mathbf{u}_1, \ldots, \mathbf{u}_m)$ is a linearly independent list of vectors of $\mathcal{V}$, it can be extended to a basis $(\mathbf{u}_1, \ldots \mathbf{u}_m, \mathbf{v}_1, \ldots, \mathbf{v}_m)$ of $\mathcal{V}$. Let $\mathcal{W} = \text{span}(\mathbf{w}_1, \ldots, \mathbf{w}_n)$. To prove $\mathcal{V} = \mathcal{U} \oplus \mathcal{W}$, need to show that $\mathcal{V} = \mathcal{U} + \mathcal{W}$ and $\mathcal{U} \cap \mathcal{W} = \{ \mathbf{0} \}$.\\
Suppose $\mathbf{v} \in \mathcal{V}$. As $\mathbf{u}_1, \ldots, \mathbf{u}_m$ spans $\mathcal{V}$, there exists scalars $a_1, \ldots, a_m, b_1, \ldots, b_n \in \mathbb{F}$ such that
\begin{align}
\mathbf{v} = a_1 \mathbf{u}_1 + \cdots + a_m \mathbf{u}_m + b_1 \mathbf{w}_1 + \cdots + b_n \mathbf{w}_n = \mathbf{u} + \mathbf{w} \nonumber
\end{align}
where $\mathbf{u} \in \mathcal{U}$, $\mathbf{w} \in \mathcal{W}$ are defined as above. Hence $\mathbf{v} \in \mathcal{U} + \mathcal{W}$, thus $\mathcal{V} = \mathcal{U} + \mathcal{W}$.\\
To show $\mathcal{U} \cap \mathcal{W} = \{ \mathbf{0} \}$, suppose $\mathbf{v} \in \mathcal{U} \cap \mathcal{W}$. Then $\exists$ scalars $a_1, \ldots, a_m, b_1, \ldots, b_n \in \mathbb{F}$. such that
\begin{align}
\mathbf{v} = a_1 \mathbf{u}_1 + \cdots + a_m \mathbf{u}_m = b_1 \mathbf{w}_1 + \cdots + b_n \mathbf{w}_n \nonumber
\end{align}
Hence
\begin{align}
a_1 \mathbf{u}_1 + \cdots + a_m \mathbf{u}_m - b_1 \mathbf{w}_1 - \cdots - b_n \mathbf{w}_n = \mathbf{0} \nonumber
\end{align}
As $(\mathbf{u}_1, \ldots, \mathbf{u}_m, \mathbf{w}_1, \ldots, \mathbf{w}_n)$ is linearly independent, this implies that $a_1 = \cdots = a_m = b_1 = \cdots = b_n = 0$, hence $\mathbf{v} = \mathbf{0}$, completing proof that $\mathcal{U} \cap \mathcal{W} = \{ \mathbf{0} \}$.
\end{proof}
\end{proposition}

\begin{theorem} Any two bases of a finite-dimensional vector space have the same length.
\begin{proof}
Suppose $\mathcal{V}$ is finite dimensional, and $\mathcal{B}_1, \mathcal{B}_2$ be any two bases of $\mathcal{V}$. Then $\mathcal{B}_1$ is linearly independent in $\mathcal{V}$, and $\mathcal{B}_2$ spans $\mathcal{V}$. Hence length of $\mathcal{B}_1$ is at most the length of $\mathcal{B}_2$. Interchanging roles of $\mathcal{B}_1$ and $\mathcal{B}_2$, the length of $\mathcal{B}_2$ is at most the length of $\mathcal{B}_1$. Hence length of $\mathcal{B}_1$ must equal length of $\mathcal{B}_2$ as desired.
\end{proof}
\end{theorem}

\begin{definition} The \hlt{dimension} of a finite-dimensional vector space is defined to be the length of any basis of the vector space, denoted as $\text{dim}(\mathcal{V})$.
\end{definition}

\begin{proposition} If $\mathcal{V}$ is finite dimensional and $\mathcal{U}$ is a subspace of $\mathcal{V}$, then $\text{dim}(\mathcal{U}) \leq \text{dim}(\mathcal{V})$.
\begin{proof} Suppose $\mathcal{V}$ is finite dimensional and $\mathcal{U}$ is a subspace of $\mathcal{V}$. Any basis of $\mathcal{U}$ is a linearly independent list of vectors in $\mathcal{V}$, thus can be extended to a basis of $\mathcal{V}$. Hence length of a basis of $\mathcal{U} \leq$ length of a basis of $\mathcal{V}$.
\end{proof}
\end{proposition}

\begin{proposition} If $\mathcal{V}$ is finite dimensional, then every spanning list of vectors in $\mathcal{V}$ with length $\text{dim}(\mathcal{V})$ is a basis of $\mathcal{V}$.
\begin{proof} Suppose $\text{dim}(\mathcal{V}) = n$ and $(\mathbf{v}_1, \ldots, \mathbf{v}_n)$ spans $\mathcal{V}$. The list $(\mathbf{v}_1, \ldots, \mathbf{v}_n)$ can be reduced to a basis of $\mathcal{V}$. However, every basis of $\mathcal{V}$ has length $n$, hence reduction must be trivial, i.e. no elements are deleted from $(\mathbf{v}_1, \ldots, \mathbf{v}_n)$. Hence, $(\mathbf{v}_1, \ldots, \mathbf{v}_n)$ is a basis of $\mathcal{V}$ as desired.
\end{proof}
\end{proposition}

\begin{proposition} If $\mathcal{V}$ is finite dimensional, then every linearly independent list of vectors in $\mathcal{V}$ with length $\text{dim}(\mathcal{V})$ is a basis of $\mathcal{V}$.
\begin{proof} Suppose $\text{dim}(\mathcal{V}) = n$ and $(\mathbf{v}_1, \ldots, \mathbf{v}_n)$ is linearly independent in $\mathcal{V}$. The list $(\mathbf{v}_1, \ldots, \mathbf{v}_n)$ can be extended to a basis of $\mathcal{V}$. However, every basis of $\mathcal{V}$ has length $n$, hence the extension must be trivial, i.e., no new elements are adjoined to $(\mathbf{v}_1, \ldots, \mathbf{v}_n)$. Hence, $(\mathbf{v}_1, \ldots, \mathbf{v}_n)$ is a basis of $\mathcal{V}$ as desired.
\end{proof}
\end{proposition}

The next theorem gives formula for dimension of sum of two subspaces of a finite-dimensional vector space.

\begin{theorem} If $\mathcal{U}_1$ and $\mathcal{U}_2$ are subspaces of a finite-dimensional vector space, then
\begin{align}
\text{dim}(\mathcal{U}_1 + \mathcal{U}_2) = \text{dim}(\mathcal{U}_1) + \text{dim}(\mathcal{U}_2) - \text{dim}(\mathcal{U}_1 \cap \mathcal{U}_2) \nonumber
\end{align}
\begin{proof} Let $(\mathbf{u}_1, \ldots, \mathbf{u}_m)$ be basis of $\mathcal{U}_1 \cap \mathcal{U}_2$, hence $\text{dim}(\mathcal{U}_1 \cap \mathcal{U}_2) = m$.\\
As $(\mathbf{u}_1, \ldots, \mathbf{u}_m)$ is a basis of $\mathcal{U}_1 \cap \mathcal{U}_2$, it is linearly independent in $\mathcal{U}_1$, hence can be extended to a basis $(\mathbf{u}_1, \ldots, \mathbf{u}_m, \mathbf{v}_1, \ldots, \mathbf{v}_j)$ of $\mathcal{U}_1$. Hence $\text{dim}(\mathcal{U}_1) = m + j$. Extend $(\mathbf{u}_1, \ldots, \mathbf{u}_m)$ to a basis of $(\mathbf{u}_1, \ldots, \mathbf{u}_m, \mathbf{w}_1, \ldots, \mathbf{w}_k)$ of $\mathcal{U}_2$, hence $\text{dim}(\mathcal{U}_2) = m + k$.\\
Note that $\text{span}(\mathbf{u}_1, \ldots, \mathbf{u}_m, \mathbf{v}_1, \ldots, \mathbf{v}_j, \mathbf{w}_1, \ldots, \mathbf{w}_k)$ contains $\mathcal{U}_1$ and $\mathcal{U}_2$, hence contains $\mathcal{U}_1 + \mathcal{U}_2$. To show that this list is a basis of $\mathcal{U}_1 + \mathcal{U}_2$, need to show that it is linearly independent. Suppose
\begin{align}
a_1 \mathbf{u}_1 + \cdots + a_m \mathbf{u}_m + b_1 \mathbf{v}_1 + \cdots + b_j \mathbf{v}_j + c_1 \mathbf{w}_1 + \cdots + c_k \mathbf{w}_k = \mathbf{0} \nonumber
\end{align}
where $a_i, b_i, c_i \in \mathbb{F}$. Rewriting the equation,
\begin{align}
c_1 \mathbf{w}_1 + \cdots + c_k \mathbf{w}_k = - a_1 \mathbf{u}_1 - \cdots - a_m \mathbf{u}_m - b_1 \mathbf{v}_1 - \cdots - b_j \mathbf{v}_j \nonumber
\end{align}
which shows $c_1 \mathbf{w}_1 + \cdots + c_k \mathbf{w}_k \in \mathcal{U}_1$. All $\mathbf{w}$ are in $\mathcal{U}_2$, hence this implies $c_1 \mathbf{w}_1 + \cdots + c_k \mathbf{w}_k \in \mathcal{U}_1 \cap \mathcal{U}_2$. As $(\mathbf{u}_1, \ldots, \mathbf{u}_m)$ is a basis of $\mathcal{U}_1 \cap \mathcal{U}_2$, write this as
\begin{align}
c_1 \mathbf{w}_1 + \cdots + c_k \mathbf{w}_k = d_1 \mathbf{u}_1 + \cdots + d_m \mathbf{u}_m \nonumber
\end{align}
for some $d_i \in \mathbb{F}$. As $(\mathbf{u}_1, \ldots, \mathbf{u}_m, \mathbf{w}_1, \ldots, \mathbf{w}_k)$ is linearly independent, the last equation implies all $c_i$ and $d_i$ equal $0$. Hence the original equation becomes
\begin{align}
a_1 \mathbf{u}_1 + \cdots + a_m \mathbf{u}_m + b_1 \mathbf{v}_1 + \cdots + b_j \mathbf{v}_j = \mathbf{0} \nonumber
\end{align}
This implies all $a_i$ and $b_i$ are $0$ as the list $(\mathbf{u}_1, \ldots, \mathbf{u}_m, \mathbf{w}_1, \ldots, \mathbf{w}_k)$ is linearly independent.\\
Hence $(\mathbf{u}_1, \ldots, \mathbf{u}_m, \mathbf{v}_1, \ldots, \mathbf{v}_j, \mathbf{w}_1, \ldots, \mathbf{w}_k)$ is a basis of $\mathcal{U}_1 + \mathcal{U}_2$, and thus
\begin{align}
\text{dim}(\mathcal{U}_1 + \mathcal{U}_2) &= m + j + k \nonumber \\
&= (m+j) + (m+k) - m \nonumber \\
&= \text{dim}(\mathcal{U}_1) + \text{dim}(\mathcal{U}_2) - \text{dim}(\mathcal{U}_1 \cap \mathcal{U}_2) \nonumber
\end{align}
\end{proof}
\end{theorem}

Note that dimension meshes well with direct sums, which is analogous to disjoint union.

\begin{proposition} Suppose $\mathcal{V}$ is finite dimensional and $\mathcal{U}_1, \ldots, \mathcal{U}_m$ are subspaces of $\mathcal{V}$ such that
\begin{align}
\mathcal{V} &= \mathcal{U}_1 + \cdots + \mathcal{U}_m \nonumber \\
\text{dim}(\mathcal{V}) &= \text{dim}(\mathcal{U}_1) + \cdots + \text{dim}(\mathcal{U}_m) \nonumber
\end{align}
Then $\mathcal{V} = \mathcal{U}_1 \oplus \cdots \oplus \mathcal{U}_m$.
\begin{proof}
Choose a basis for each $\mathcal{U}_j$. Put these bases together in a list which spans $\mathcal{V}$ and has length $\text{dim}(\mathcal{V})$. Hence this list is a basis of $\mathcal{V}$, and it is linearly independent. Suppose $\mathbf{u}_1 \in \mathcal{U}_1, \ldots, \mathbf{u}_m \in \mathcal{U}_m$ such that
\begin{align}
\mathbf{0} = \mathbf{u}_1 + \cdots + \mathbf{u}_m \nonumber
\end{align}
Write each $\mathbf{u}_j$ as a linear combination of basis vectors of $\mathcal{U}_j$. Substitute these linear combinations into equation above, hence $\mathbf{0}$ is written as a linear combination of basis of $\mathcal{V}$ constructed above. Hence all scalars used in this linear combination must be $0$. Thus each $\mathbf{u}_j = \mathbf{0}$, which proves $\mathcal{V} = \mathcal{U}_1 \oplus \cdots \oplus \mathcal{U}_m$.
\end{proof}
\end{proposition}

\subsubsection{Linear Maps}

\begin{definition} A \hlt{Linear Map} from $\mathcal{V}$ to $\mathcal{W}$ is a function $T: \mathcal{V} \rightarrow \mathcal{W}$ with following properties:
\begin{enumerate}[label=\roman*.]
\setlength{\itemsep}{0pt}
\item \hlt{Additivity}: $T(\mathbf{u} + \mathbf{v}) = T\mathbf{u} + T\mathbf{v}$ for all $\mathbf{u}, \mathbf{v} \in \mathcal{V}$
\item \hlt{Homogeneity}: $T(a\mathbf{v}) = a(T\mathbf{v})$ for all $a \in \mathbb{F}$ and all $\mathbf{v} \in \mathcal{V}$.
\end{enumerate}
The set of all linear maps from $\mathcal{V}$ to $\mathcal{W}$ is denoted $\mathcal{L}(\mathcal{V}, \mathcal{W})$.\\
The linear map is \hlt{injective} if whenever $\mathbf{u}, \mathbf{v} \in \mathcal{V}$ and $T\mathbf{u} = T\mathbf{v}$, that $\mathbf{u} = \mathbf{v}$.\\
The linear map is \hlt{surjective} if its range equals $\mathcal{W}$.
\end{definition}

\begin{example} Examples of linear maps include
\begin{enumerate}[label=\roman*.]
\setlength{\itemsep}{0pt}
\item Zero Map: $0 \in \mathcal{L}(\mathcal{V}, \mathcal{W})$ is defined by $0 \mathbf{v} = 0$.
\item Identity Map: $I \in \mathcal{L}(\mathcal{V}, \mathcal{V})$ is defined by $I \mathbf{v} = \mathbf{v}$
\item Differentiation: $T \in \mathcal{L}(\mathcal{P}(\mathbf{R}), \mathcal{P}(\mathbf{R}))$ is defined by $T \mathbf{p} = \mathbf{p}'$
\item Integration: $T \in \mathcal{L}(\mathcal{P}(\mathbf{R}), \mathbf{R})$ is defined by $T \mathbf{p} = \int_0^1 \mathbf{p}(\mathbf{x}) \ d \mathbf{x}$
\item Backward Shift: Let $\mathbb{F}^{\infty}$ be vector space of all sequences of elements of $\mathbb{F}$. Let $T \in \mathcal{L}(\mathbb{F}^{\infty}, \mathbb{F}^{\infty})$. The backward shift map is then $T(x_1, x_2, x_3, \ldots) = (x_2, x_3, \ldots)$
\item From $\mathbb{F}^{n}$ to $\mathbb{F}^{m}$: $T \in \mathcal{L}(\mathbb{F}^{n}, \mathbb{F}^{m})$ defined by $T(x_1, \ldots, x_n) = (a_{1,1} x_1 + \cdots + a_{1,n} x_n, \ldots, a_{m,1} x_1 + \cdots + a_{m,n} x_n)$, where $a_{j,k} \in \mathbb{F}$, $j=1, \ldots, m$ and $k=1, \ldots, n$.
\end{enumerate}
\end{example}

\begin{remark} \hlt{Properties of Product of Linear Maps}\\
Let $T \in \mathcal{L}(\mathcal{U}, \mathcal{V})$, and $S \in \mathcal{L}(\mathcal{V}, \mathcal{W})$. The \hlt{product} of $S$ and $T$ is $ST$.
\begin{enumerate}[label=\roman*.]
\setlength{\itemsep}{0pt}
\item Associativity: $(T_1 T_2)T_3 = T_1 (T_2 T_3)$
\item Identity: $TI = T$ (identity map on $\mathcal{V}$) and $IT = T$ (identity map on $\mathcal{W}$)
\item Distributive: $(S_1 + S_2)T = S_1 T + S_2 T$ and $S(T_1 + T_2) = ST_1 + ST_2$ whenever $T, T_1, T_2 \in \mathcal{L}(\mathcal{U}, \mathcal{V})$ and $S, S_1, S_2 \in \mathcal{L}(\mathcal{V}, \mathcal{W})$
\end{enumerate}
Multiplication of linear maps is not commutative.
\end{remark}

\begin{definition} For $T \in \mathcal{L}(\mathcal{V}, \mathcal{W})$, the \hlt{null space} of $T$, denoted $\text{null}(T)$, is the subset of $\mathcal{V}$ consisting of those vectors that maps $T$ to $0$, $\text{null}(T) = \{ \mathbf{v} \in \mathcal{V}: T \mathbf{v} = \mathbf{0} \}$.
\end{definition}

\begin{proposition} \label{prop:nullspaceissubspace} If $T \in \mathcal{L}(\mathcal{V}, \mathcal{W})$, then $\text{null}(T)$ is a subspace of $\mathcal{V}$.
\begin{proof} Suppose $T \in \mathcal{L}(\mathcal{V}, \mathcal{W})$. By additivity,
\begin{align}
T(\mathbf{0}) = T(\mathbf{0} + \mathbf{0}) = T(\mathbf{0}) + T(\mathbf{0}) \nonumber
\end{align}
which implies $T(\mathbf{0}) = \mathbf{0}$. Hence $\mathbf{0} \in \text{null}(T)$.\\
If $\mathbf{u}, \mathbf{v} \in \text{null}(T)$, then
\begin{align}
T(\mathbf{u} + \mathbf{v}) = T\mathbf{u} + T\mathbf{v} = \mathbf{0} + \mathbf{0} = \mathbf{0} \nonumber
\end{align}
Hence $\mathbf{u} + \mathbf{v} \in \text{null}(T)$, and $\text{null}(T)$ is closed under addition.\\
If $\mathbf{u} \in \text{null}(T)$ and $a \in \mathbb{F}$, then
\begin{align}
T(a \mathbf{u}) = a T \mathbf{u} = a\mathbf{0} = \mathbf{0} \nonumber
\end{align}
Hence $a \mathbf{u} \in \text{null}(T)$, and $\text{null}(T)$ is closed under scalar multiplication.\\
As $\text{null}(T)$ contains $\mathbf{0}$ and is closed under addition and scalar multiplication, $\text{null}(T)$ is subspace of $\mathcal{V}$.
\end{proof}
\end{proposition}

\begin{proposition} Let $T \in \mathcal{L}(\mathcal{V}, \mathcal{W})$. Then $T$ is injective if and only if $\text{null}(T) = \{ \mathbf{0} \}$.
\begin{proof} ($\Rightarrow$) Suppose $T$ is injective. By Proposition \ref{prop:nullspaceissubspace}, $\{ \mathbf{0} \} \subset \text{null}(T)$. Suppose $\mathbf{v} \in \text{null}(T)$. Then $T(\mathbf{v}) = \mathbf{0} = T(\mathbf{0})$. As $T$ is injective, equation above implies $\mathbf{v} = \mathbf{0}$. Thus $\text{null}(T) = \{ \mathbf{0} \}$.\\
($\Leftarrow$) Suppose $\text{null}(T) = \{ \mathbf{0} \}$, $\mathbf{u}, \mathbf{v} \in \mathcal{V}$ and $T \mathbf{u} = T \mathbf{v}$. Then $\mathbf{0} = T \mathbf{u} - T\mathbf{v} = T(\mathbf{u} - \mathbf{v})$. Hence $\mathbf{u} - \mathbf{v}$ is in $\text{null}(T)$, which equals $\{ \mathbf{0} \}$. Thus $\mathbf{u} - \mathbf{v} = \mathbf{0}$, which implies $\mathbf{u} = \mathbf{v}$. Hence $T$ is injective.
\end{proof}
\end{proposition}

\begin{definition} For $T \in \mathcal{L}(\mathcal{V}, \mathcal{W})$, the \hlt{range} of $T$, denoted $\text{range}(T)$, is subset of $\mathcal{W}$ consisting of vectors that are of from $T \mathbf{v}$ for some $\mathbf{v} \in \mathcal{V}$, where $\text{range}(T) = \{ T \mathbf{v} : \mathbf{v} \in \mathcal{V} \}$.
\end{definition}

\begin{proposition} If $T \in \mathcal{L}(\mathcal{V}, \mathcal{W})$, then $\text{range}(T)$ is a subspace of $\mathcal{W}$.
\begin{proof} Suppose $T \in \mathcal{L}(\mathcal{V}, \mathcal{W})$, then $T(\mathbf{0}) = \mathbf{0}$ by Proposition \ref{prop:nullspaceissubspace}, implying $\mathbf{0} \in \text{range}(T)$.\\
If $\mathbf{w}_1, \mathbf{w}_2 \in \text{range}(T)$, then $\exists \ \mathbf{v}_1, \mathbf{v}_2 \in \mathcal{V}$ such that $T \mathbf{v}_1 = \mathbf{w}_1$ and $T \mathbf{v}_2 = \mathbf{w}_2$. Hence
\begin{align}
T(\mathbf{v}_1 + \mathbf{v}_2) = T\mathbf{v}_1 + T\mathbf{v}_2 = \mathbf{w}_1 + \mathbf{w}_2 \nonumber
\end{align}
thus $\mathbf{w}_1 + \mathbf{w}_2 \in \text{range}(T)$. Hence $\text{range}(T)$ is closed under addition.\\
If $\mathbf{w} \in \text{range}(T)$ and $a \in \mathbb{F}$, then $\exists \ \mathbf{v} \in \mathcal{V}$ such that $T \mathbf{v} = \mathbf{w}$. Hence $T(a \mathbf{v}) = a T \mathbf{v} = a \mathbf{w}$, and $a \mathbf{w} \in \text{range}(T)$. Thus $\text{range}(T)$ is closed under scalar multiplication.\\
As $\text{range}(T)$ contains $\mathbf{0}$ and is closed under addition and scalar multiplication, $\text{range}(T)$ is a subspace of $\mathcal{W}$.
\end{proof}
\end{proposition}

The next theorem states that the dimension of the null space plus dimension of range of linear map on finite-dimensional vector space equals dimension of the domain.


\begin{theorem} If $\mathcal{V}$ is finite-dimensional and $T \in \mathcal{L}(\mathcal{V}, \mathcal{W})$, then $\text{range}(T)$ is a finite dimensional subspace of $\mathcal{W}$, and $\text{dim}(\mathcal{V}) = \text{dim}(\text{null}(T)) + \text{dim}(\text{range}(T))$.
\begin{proof} Suppose $\mathcal{V}$ is a finite-dimensional vector space and $T \in \mathcal{L}(\mathcal{V}, \mathcal{W})$.\\
Let $(\mathbf{u}_1, \ldots, \mathbf{u}_m)$ be a basis of $\text{null}(T)$, thus $\text{dim}(\text{null}(T)) = m$. The linearly independent list $(\mathbf{u}_1, \ldots, \mathbf{u}_m)$ can be extended to a basis $(\mathbf{u}_1, \ldots, \mathbf{u}_m, \mathbf{w}_1, \ldots, \mathbf{w}_n)$ of $\mathcal{V}$. Hence $\text{dim}(\mathcal{V}) = m + n$. Need to show $\text{range}(T)$ is finite dimensional and $\text{dim}(\text{range}(T)) = n$, by proving that $(T \mathbf{w}_1, \ldots, T \mathbf{w}_n)$ is a basis of $\text{range}(T)$.\\
Let $\mathbf{v} \in \mathcal{V}$. As $(\mathbf{u}_1, \ldots, \mathbf{u}_m, \mathbf{w}_1, \ldots, \mathbf{w}_n)$ spans $\mathcal{V}$, then
\begin{align}
\mathbf{v} = a_1 \mathbf{u}_1 + \cdots + a_m \mathbf{u}_m + b_1 \mathbf{w}_1 + \cdots + b_n \mathbf{w}_n \nonumber
\end{align}
where $a, b \in \mathbb{F}$. Applying $T$ to both sides of equation,
\begin{align}
T\mathbf{v} = b_1 T \mathbf{w}_1 + \cdots + b_n T \mathbf{w}_n \nonumber
\end{align}
Note $T \mathbf{u}_j$ disappeared as each $\mathbf{u}_j \in \text{null}(T)$. Last equation implies $(T \mathbf{w}_1, \ldots T \mathbf{w}_n)$ spans $\text{range}(T)$, hence that $\text{range}(T)$ is finite dimensional.\\
To show $(T \mathbf{w}_1, \ldots T \mathbf{w}_n)$ is linearly independent, suppose $c_1, \ldots, c_n \in \mathbb{F}$, and
\begin{align}
c_1 T \mathbf{w}_1 + \cdots + c_n T \mathbf{w}_n &= \mathbf{0} \nonumber \\
T(c_1 \mathbf{w}_1 + \cdots + c_n \mathbf{w}_n) &= \mathbf{0} \nonumber \\
c_1 \mathbf{w}_1 + \cdots + c_n \mathbf{w}_n &\in \text{null}(T) \nonumber
\end{align}
As $(\mathbf{u}_1, \ldots, \mathbf{u}_m)$ spans $\text{null}(T)$, then
\begin{align}
c_1 \mathbf{w}_1 + \cdots + c_n \mathbf{w}_n = d_1 \mathbf{u}_1 + \cdots + d_m \mathbf{u}_m \nonumber
\end{align}
where $d_i \in \mathbb{F}$. This implies all $c$ and $d$ are $0$, as $(\mathbf{u}_1, \ldots, \mathbf{u}_m, \mathbf{w}_1, \ldots, \mathbf{w}_n)$ are linearly independent.\\
Hence $(T \mathbf{w}_1, \ldots, T \mathbf{w}_n)$ is linearly independent and is a basis for $\text{range}(T)$.
\end{proof}
\end{theorem}

No linear map from finite-dimensional vector space to a smaller vector space can be injective.

\begin{corollary} If $\mathcal{V}$ and $\mathcal{W}$ are finite-dimensional vector spaces such that $\text{dim}(\mathcal{V}) > \text{dim}(\mathcal{W})$, then no linear map from $\mathcal{V}$ to $\mathcal{W}$ is injective.
\begin{proof} Suppose $\mathcal{V}$ and $\mathcal{W}$ are finite-dimensional vector spaces such that $\text{dim}(\mathcal{V}) > \text{dim}(\mathcal{W})$. Let $T \in \mathcal{L}(\mathcal{V}, \mathcal{W})$,
\begin{align}
\text{dim}(\text{null}(T)) &= \text{dim}(\mathcal{V}) - \text{dim}(\text{range}(T)) \nonumber \\
&\geq \text{dim}(\mathcal{V}) - \text{dim}(\mathcal{W}) \geq 0 \nonumber
\end{align}
As $\text{dim}(\text{null}(T)) \geq 0$, $\text{null}(T)$ must contain vectors other than $0$, hence $T$ is not injective.
\end{proof}
\end{corollary}

No linear map from finite-dimensional vector space to a bigger vector space can be surjective.

\begin{corollary} If $\mathcal{V}$ and $\mathcal{W}$ are finite-dimensional vector spaces such that $\text{dim}(\mathcal{V}) < \text{dim}(\mathcal{W})$, then no linear map from $\mathcal{V}$ to $\mathcal{W}$ is surjective.
\begin{proof} Suppose $\mathcal{V}$ and $\mathcal{W}$ are finite-dimensional vector spaces such that $\text{dim}(\mathcal{V}) < \text{dim}(\mathcal{W})$. Let $T \in \mathcal{L}(\mathcal{V}, \mathcal{W})$,
\begin{align}
\text{dim}(\text{range}(T)) &= \text{dim}(\mathcal{V}) - \text{dim}(\text{null}(T)) \nonumber \\
&\leq \text{dim}(\mathcal{V}) < \text{dim}(\mathcal{W}) \nonumber
\end{align}
As $\text{dim}(\text{range}(T)) < \text{dim}(\mathcal{W)}$, $\text{range}(T)$ cannot equal $\mathcal{W}$, hence $T$ is not surjective.
\end{proof}
\end{corollary}

\begin{definition} A $m \times n$ \hlt{matrix} $\mathcal{M}$ has $m$ rows and $n$ columns,
\begin{align}
\begin{bmatrix}
a_{1,1} & \cdots & a_{1,n} \\
\vdots & \ddots & \vdots \\
a_{m,1} & \cdots & a_{m,n}
\end{bmatrix} \nonumber
\end{align}
\end{definition}

\begin{proposition} Suppose $T \in \mathcal{L}(\mathcal{V}, \mathcal{W})$, and $(\mathbf{v}_1, \ldots, \mathbf{v}_n)$ is a basis of $\mathcal{V}$, and $(\mathbf{w}_1, \ldots, \mathbf{w}_m)$ is a basis of $\mathcal{W}$. Then $\mathcal{M}(T\mathbf{v}) = \mathcal{M}(T)\mathcal{M}(\mathbf{v})$ for every $\mathbf{v} \in \mathcal{V}$.
\begin{proof} Let
\begin{align}
\begin{bmatrix}
a_{1,1} & \cdots & a_{1,n} \\
\vdots & \ddots & \vdots \\
a_{m,1} & \cdots & a_{m,n}
\end{bmatrix} \nonumber
\end{align}
Note that
\begin{align}
T\mathbf{v}_k = \sum\limits_{j=1}^{m} a_{j,k} \mathbf{w}_{j} \nonumber
\end{align}
for each $k$. Let $\mathbf{v}$ be an arbitrary vector in $\mathcal{V}$, then $\mathcal{M}(\mathbf{v})$ is given by
\begin{align}
T\mathbf{v} &= b_1 T \mathbf{v}_1 + \cdots + b_n T \mathbf{v}_n \nonumber \\
&= b_1 \sum\limits_{j=1}^{m} a_{j,1} \mathbf{w}_j + \cdots + b_n \sum\limits_{j=1}^{m} a_{j,n} \mathbf{w}_j \nonumber \\
&= \sum\limits_{j=1}(a_{j,1} b_1 + \cdots + a_{j,n} b_n) \mathbf{w}_j \nonumber
\end{align}
The equation shows that $\mathcal{M}(T \mathbf{v})$ is a $m \times 1$ matrix of vector $T \mathbf{v}$ with respect to basis $(\mathbf{w}_1, \ldots, \mathbf{w}_m)$, given by
\begin{align}
\mathcal{M}(T\mathbf{v}) = 
\begin{bmatrix}
a_{1,1} b_1 + \cdots + a_{1,n} b_n \\
\vdots \\
a_{m,1} b_1 + \cdots + a_{m,n} b_n
\end{bmatrix} \nonumber
\end{align}
Hence $\mathcal{M}(T\mathbf{v}) = \mathcal{M}(T)\mathcal{M}(\mathbf{v})$.
\end{proof}
\end{proposition}

\begin{definition} A linear map $T \in \mathcal{L}(\mathcal{V}, \mathcal{W})$ is \hlt{invertible} if $\exists$ linear map $S \in \mathcal{L}(\mathcal{W}, \mathcal{V})$ such that $ST$ equals identity map on $\mathcal{V}$ and $TS$ equals identity map on $\mathcal{W}$.\\
A linear map $S \in \mathcal{L}(\mathcal{W}, \mathcal{V})$ satisfying $ST = I$ and $TS=I$ is the \hlt{inverse} of $T$.
\end{definition}

\begin{proposition} A linear map is invertible if and only if it is injective and surjective.
\begin{proof} Let $T \in \mathcal{L}(\mathcal{V}, \mathcal{W})$.\\
($\Rightarrow$) Suppose $T$ is invertible, $\mathbf{u}, \mathbf{v} \in \mathcal{V}$, and $T \mathbf{u} = T\mathbf{v}$. Then
\begin{align}
\mathbf{u} = T^{-1}(T \mathbf{u}) = T^{-1}T(\mathbf{v}) = \mathbf{v} \nonumber
\end{align}
hence $\mathbf{u} = \mathbf{v}$. Thus $T$ is injective.\\
To show $T$ is surjective, let $\mathbf{w} \in \mathcal{W}$, then $\mathbf{w} = T (T^{-1} \mathbf{w})$, thus $\mathbf{w}$ is in range of $T$. Hence $\text{range}(T) = W$.\\
Suppose $T$ is injective and surjective, to prove $T$ is invertible. For each $\mathbf{w} \in \mathcal{W}$, define $S \mathbf{w}$ to be unique element of $\mathcal{V}$ such that $T(S \mathbf{w}) = \mathbf{w}$, which shows $TS$ equals identity map on $\mathcal{W}$. To prove $ST$ equals identity map on $\mathcal{V}$, let $\mathbf{v} \in \mathcal{V}$, then
\begin{align}
T(ST \mathbf{v}) = (TS)(T\mathbf{v}) = I(T\mathbf{v}) = T\mathbf{v} \nonumber
\end{align}
which implies $ST \mathbf{v} = \mathbf{v}$ (as $T$ is injective), thus $ST$ equals identity map on $\mathcal{V}$. Let $\mathbf{w}_1, \mathbf{w}_2 \in \mathcal{W}$, then
\begin{align}
T(S\mathbf{w}_1 + S\mathbf{w}_2) = T(S\mathbf{w}_1) + T(S\mathbf{w}_2) = \mathbf{w}_1 + \mathbf{w}_2 \nonumber
\end{align}
Thus $S\mathbf{w}_1 + S\mathbf{w}_2$ is unique element of $\mathcal{V}$ that maps $T$ to $\mathbf{w}_1 + \mathbf{w}_2$. This implies $S$ satisfies the additive property requires for linearity.\\
For proof of homogeneity, let $\mathbf{w} \in \mathcal{W}$ and $a \in \mathbb{F}$, then
\begin{align}
T(a S \mathbf{w}) = a T(S \mathbf{w}) = a \mathbf{w} \nonumber
\end{align}
Thus $a S \mathbf{w}$ is the unique element of $\mathcal{V}$ that $T$ maps to $a \mathbf{w}$. By definition of $S$, this implies $S(a\mathbf{w}) = aS\mathbf{w}$, hence $S$ is linear, as desired.
\end{proof}
\end{proposition}

\begin{definition} Two vector spaces are \hlt{isomorphic} if there is an invertible linear map from one vector space onto the other. Two isomorphic spaces have the same properties.
\end{definition}

\begin{theorem} \label{thm:isomorphicfindimvectorspace} Two finite-dimensional vector spaces are isomorphic if and only if they have the same $\text{dim}$.
\begin{proof} ($\Rightarrow$) Suppose $\mathcal{V}$ and $\mathcal{W}$ are isomorphic finite-dimensional vector spaces.\\
$\exists$ invertible linear map $T \in \mathcal{L}(\mathcal{V}, \mathcal{W})$. As $T$ is invertible, $\text{null}(T) = \{ \mathbf{0} \}$, and $\text{range}(T) = \mathcal{W}$, thus $\text{dim}(\text{null}(T)) = 0$ and $\text{dim}(\text{range}(T)) = \text{dim}(\mathcal{W})$.
\begin{align}
\text{dim}(\mathcal{V}) &= \text{dim}(\text{null}(T)) + \text{dim}(\text{range}(T)) \nonumber \\
&= \text{dim}(\mathcal{W}) \nonumber
\end{align}
($\Leftarrow$) Suppose $\mathcal{V}$ and $\mathcal{W}$ are finite-dimensional vector spaces with the same dimension. let $(\mathbf{v}_1, \ldots, \mathbf{v}_n)$ be basis of $\mathcal{V}$ and $(\mathbf{w}_1, \ldots, \mathbf{w}_n)$ be basis of $\mathcal{W}$. Let $T \in \mathcal{L}(\mathcal{V}, \mathcal{W})$ defined by
\begin{align}
T(a_1 \mathbf{v}_1 + \cdots + a_n \mathbf{v}_n) = a_1 \mathbf{w}_1 + \cdots + a_n \mathbf{w}_n \nonumber
\end{align}
Then $T$ is surjective as $(\mathbf{w}_1, \ldots, \mathbf{w}_n)$ spans $\mathcal{W}$, and $T$ is injective as $(\mathbf{w}_1, \ldots, \mathbf{w}_n)$ is linearly independent.\\
As $T$ is injective and surjective, it is invertible, hence $\mathcal{V}$ and $\mathcal{W}$ are isomorphic.
\end{proof}
\end{theorem}

The next theorem implies that every finite-dimensional vector space is isomorphic to some $\mathbb{F}^n$.

\begin{theorem} \label{thm:invertiblelinmap} Suppose $(\mathbf{v}_1, \ldots, \mathbf{v}_n)$ is a basis of $\mathcal{V}$ and $(\mathbf{w}_1, \cdots, \mathbf{w}_m)$ is a basis of $\mathcal{W}$. Then $\mathcal{M}$ is an invertible linear map between $\mathcal{L}(\mathcal{V}, \mathcal{W}) $ and $\text{Mat}(m,n,\mathbb{F})$,
\begin{proof} Note that $\mathcal{M}$ is linear, hence to prove that $\mathcal{M}$ is injective and surjective.\\
If $T \in \mathcal{L}(\mathcal{V}, \mathcal{W})$ and $\mathcal{M}(T) = 0$, then $T\mathbf{v}_k = \mathbf{0}$ for $k = 1, \ldots, n$, As $(\mathbf{v}_1, \ldots, \mathbf{v}_n)$ is a basis of $\mathcal{V}$, this implies $T=0$, hence $\mathcal{M}$ is injective.\\
To show $\mathcal{M}$ is surjective, let
\begin{align}
A =
\begin{bmatrix}
a_{1,1} & \cdots & a_{1,n} \\
\vdots & \ddots & \vdots \\
a_{m,1} & \cdots & a_{m,n}
\end{bmatrix} \nonumber
\end{align}
be a matrix in $\text{Mat}(m,n,\mathbb{F})$. Let $T \in \mathcal{L}(\mathcal{V}, \mathcal{W})$ such that
\begin{align}
T \mathbf{v}_k = \sum\limits_{j=1}^m a_{j,k} \mathbf{w}_k \nonumber
\end{align}
for $k=1, \ldots, n$. Note $\mathcal{M}(T) = A$, so range of $\mathcal{M}$ equals $\text{Mat}(m,n,\mathbb{F})$ as desired.
\end{proof}
\end{theorem}

\begin{proposition} If $\mathcal{V}$ and $\mathcal{W}$ are finite-dimensional, then $\mathcal{L}(\mathcal{V}, \mathcal{W})$ is finite-dimensional, and
\begin{align}
\text{dim}\mathcal{L}(\mathcal{V}, \mathcal{W}) = (\text{dim}(\mathcal{V}))(\text{dim}(\mathcal{W})) \nonumber
\end{align}
\begin{proof} This follows from $\text{dim}(\text{Mat}(m,n,\mathbb{F})) = mn$ and Theorem \ref{thm:isomorphicfindimvectorspace} and Theorem \ref{thm:invertiblelinmap}.
\end{proof}
\end{proposition}

\begin{definition} A linear map from a vector space to itself is an \hlt{operator}, denoted $\mathcal{L}(\mathcal{V})$.
\end{definition}

\begin{theorem} Suppose $\mathcal{V}$ is finite-dimensional. If $T \in \mathcal{L}(\mathcal{V})$, then the following are equivalent:
\begin{enumerate}[label=\roman*.]
\setlength{\itemsep}{0pt}
\item $T$ is invertible.
\item $T$ is injective.
\item $T$ is surjective.
\end{enumerate}
\begin{proof} Suppose $T \in \mathcal{L}(\mathcal{V})$, then $(i) \Rightarrow (ii)$.\\
Suppose $(ii)$ holds so that $T$ is injective. Thus $\text{null}(T) = \{ \mathbf{0} \}$. Then
\begin{align}
\text{dim}(\text{range}(T)) &= \text{dim}(\mathcal{V}) - \text{dim}(\text{null}(T)) \nonumber \\
&= \text{dim}(\mathcal{V}) \nonumber
\end{align}
which implies range of $T$ equals $\mathcal{V}$, hence $T$ is surjective, and $(ii) \Rightarrow (iii)$.\\
Suppose $(iii)$ holds, so $T$ is surjective. Thus $\text{range}(T) = \mathcal{V}$. Then
\begin{align}
\text{dim}(\text{null}(T))  = \text{dim}(\mathcal{V}) - \text{dim}(\text{range}(T)) = 0 \nonumber 
\end{align}
which implies $\text{null}(T)$ equals $\{ \mathbf{0} \}$, hence $T$ is injective and thus invertible. Hence $(iii) \Rightarrow (i)$.
\end{proof}
\end{theorem}

\subsubsection{Eigenvalues and Eigenvectors}

\begin{definition} Let $T \in \mathcal{L}(\mathcal{V})$, and $\exists$ direct sum decomposition $\mathcal{V} = \mathcal{U}_1 \oplus \cdots \oplus \mathcal{U}_m$ where each $\mathcal{U}_j$ is proper subspace of $\mathcal{V}$. The \hlt{restriction} of $T$ to smaller domain $\mathcal{U}_j$ is denoted $T|_{\mathcal{U}_j}$
\end{definition}

\begin{definition} For $T \in \mathcal{L}(\mathcal{V})$ and $\mathcal{U}$ a subspace of $\mathcal{V}$, $\mathcal{U}$ is invariant under $T$ if $\mathbf{u} \in \mathcal{U}$ implies $T \mathbf{u} \in \mathcal{U}$ ($\mathcal{U}$ is invariant under $T$ if $T|_{\mathcal{U}}$ is an operator under $\mathcal{U}$.
\end{definition}

\begin{definition} A scalar $\lambda \in \mathbb{F}$ is an \hlt{eigenvalue} of $T \in \mathcal{L}(\mathcal{V})$ if $\exists$ a nonzero vector $\mathbf{u} \in \mathcal{V}$ such that $T \mathbf{u} = \lambda \mathbf{u}$.\\
A vector $\mathbf{u} \in \mathcal{V}$ is an \hlt{eigenvector} of $T$ if $T \mathbf{u} = \lambda \mathbf{u}$.
\end{definition}

Note that the set of eigenvectors of $T$ corresponding to $\lambda$ equals $\text{null}(T - \lambda I)$.

\begin{theorem} \label{thm:disctingeigenvectorsindependence} Let $T \in \mathcal{L}(\mathcal{V})$, and $\lambda_1, \ldots, \lambda_m$ distinct eigenvalues of $T$, $\mathbf{v}_1, \ldots, \mathbf{v}_m$ are corresponding nonzero eigenvectors, then $(\mathbf{v}_1, \ldots, \mathbf{v}_m)$ is linearly independent.
\begin{proof}
Suppose $(\mathbf{v}_1, \ldots, \mathbf{v}_m)$ is linearly dependent. Let $k$ be smallest positive integer such that
\begin{align}
\mathbf{v}_k \in \text{span}(\mathbf{v}_1, \ldots, \mathbf{v}_{k-1}) \nonumber
\end{align}
The existence of $k$ with this property follows from Lemma \ref{lem:lineardependencelemma}. Hence $\exists \ a_1, \ldots, a_{k-1} \in \mathbb{F}$ such that
\begin{align}
\mathbf{v}_k = a_1 \mathbf{v}_1 + \cdots + a_{k-1} \mathbf{v}_{k-1} \nonumber
\end{align}
Applying $T$ to both sides of equation gives
\begin{align}
\lambda_k \mathbf{v}_k = a_1 \lambda_1 \mathbf{v}_1 + \cdots + a_{k-1} \lambda_{k-1} \mathbf{v}_{k-1} \nonumber
\end{align}
Multiply both sides of second equation by $\lambda_k$ then subtracting above gives
\begin{align}
0 = a_1(\lambda_k - \lambda_1)\mathbf{v}_1 + \cdots + a_{k-1}(\lambda_k - \lambda_{k-1}) \mathbf{v}_{k-1} \nonumber
\end{align}
As $k$ is chosen to be smallest positive integer satisfying first equation, $(\mathbf{v}_1, \ldots, \mathbf{v}_{k-1})$ is linearly dependent, hence equation above implies that all $a$ are $0$, which also means $\mathbf{v}_k = \mathbf{0}$, contracting hypothesis that all $\mathbf{v}$ are nonzero. Hence assumption that $(\mathbf{v}_1, \ldots, \mathbf{v}_m)$ is linearly dependent must be false.
\end{proof}
\end{theorem}

\begin{corollary} Each operator on $\mathcal{V}$ has at most $\text{dim}(\mathcal{V})$ distinct eigenvalues.
\begin{proof} Let $T \in \mathcal{L}(\mathcal{V})$, and $\lambda_1, \ldots, \lambda_m$ distinct eigenvalues of $T$. Let $\mathbf{v}_1, \ldots, \mathbf{v}_m$ be corresponding nonzero eigenvectors. Theorem \ref{thm:disctingeigenvectorsindependence} implies $(\mathbf{v}_1, \ldots, \mathbf{v}_m)$ is linearly independent, hence $m \leq \text{dim}(\mathcal{V})$.
\end{proof}
\end{corollary}

\begin{theorem} Every operator on a finite-dimensional, nonzero complex vector space has an eigenvalue.
\begin{proof} Suppose $\mathcal{V}$ is a complex vector space with $\text{dim}(\mathcal{V})=n>0$ and $T \in \mathcal{L}(\mathcal{V})$. Choose $\mathbf{v} \in \mathcal{V}$ with $\mathbf{v} \neq \mathbf{0}$, then $(\mathbf{v}, T\mathbf{v}, T^2 \mathbf{v}, \ldots, T^n \mathbf{v})$ cannot be linearly independent as $\mathcal{V}$ has dimension $n$ and $n+1$ vectors. Thus $\exists \ a_0, \ldots, a_n \in \C$ not all $0$ such that
\begin{align}
0 = a_0 \mathbf{v} + a_1 T \mathbf{v} + \cdots + a_n T^n \mathbf{v} \nonumber
\end{align}
Let $m$ be largest index such that $a_m \neq 0$. As $\mathbf{v} \neq \mathbf{0}$, the coefficients $a_1, \ldots, a_m$ cannot all be $0$, hence $0 < m \leq n$. Let $a$ e coefficients of a polynomial in factored form
\begin{align}
a_0 + a_1 z + \cdots + a_n z^n = c(z - \lambda_1) \cdots (z - \lambda_m) \nonumber
\end{align}
where $c$ is nonzero complex number, each $\lambda_j \in \C$, and equation holds for all $z \in \C$, then
\begin{align}
\mathbf{0} &= a_0 \mathbf{v} + a_1 T \mathbf{v} + \cdots + a_n T^n \mathbf{v} \nonumber \\
&= (a_0 I + a_1 T + \cdots + a_n T^n)\mathbf{v} \nonumber \\
&= c(T - \lambda_1 I) \cdots (T - \lambda_m I)\mathbf{v} \nonumber
\end{align}
which means $T - \lambda_j I$ is not injective for at least one $j$. Hence $T$ has an eigenvalue.
\end{proof}
\end{theorem}

\begin{definition} \hlt{Upper triangular matrices} are of the form
\begin{align}
\begin{bmatrix}
\lambda_1 & & * \\
& \ddots & \\
0 & & \lambda_n
\end{bmatrix} \nonumber
\end{align}
\end{definition}

Next proposition shows useful connection between upper-triangular matrices and invariant subspaces.

\begin{proposition} Suppose $T \in \mathcal{L}(\mathcal{V})$, and $(\mathbf{v}_1, \ldots, \mathbf{v}_n)$ is a basis of $\mathcal{V}$. Following are equivalent:
\begin{enumerate}[label=\roman*.]
\setlength{\itemsep}{0pt}
\item The matrix of $T$ with respect to $(\mathbf{v}_1, \ldots, \mathbf{v}_n)$ is upper triangular.
\item $T \mathbf{v}_k \in \text{span}(\mathbf{v}_1, \ldots, \mathbf{v}_n)$ for each $k = 1, \ldots, n$
\item $\text{span}(\mathbf{v}_1, \ldots, \mathbf{v}_n)$ is invariant under $T$ for each $k=1, \ldots, n$
\end{enumerate}
\begin{proof} $(i) \Rightarrow (ii)$ follows from definitions. Also, $(iii) \Rightarrow (ii)$ as $(iii)$ implies $(ii)$.\\
$(ii) \Rightarrow (iii)$, suppose $(ii)$ holds. Fix $k \in \{1, \ldots, n \}$. From $(ii)$, we know that
\begin{align}
T \mathbf{v}_1 &\in \text{span}(\mathbf{v}_1) \subset \text{span}(\mathbf{v}_1, \ldots, \mathbf{v}_k) \nonumber \\
T \mathbf{v}_2 &\in \text{span}(\mathbf{v}_1, \mathbf{v}_2) \subset \text{span}(\mathbf{v}_1, \ldots, \mathbf{v}_k) \nonumber \\
&\vdots \nonumber \\
T \mathbf{v}_k &\in \text{span}(\mathbf{v}_1, \ldots, \mathbf{v}_k) \nonumber
\end{align}
Hence if $\mathbf{v}$ is a linear combination of $(\mathbf{v}_1, \ldots, \mathbf{v}_k)$, then
\begin{align}
T \mathbf{v} \in \text{span}(\mathbf{v}_1, \ldots, \mathbf{v}_k) \nonumber
\end{align}
Hence, $\text{span}(\mathbf{v}_1, \ldots, \mathbf{v}_k)$ is invariant under $T$.
\end{proof}
\end{proposition}

Next theorem shows that for each operator on a complex vector space, there is a basis of the vector space with respect to which the matrix of operator only has $0$ below the diagonal.

\begin{theorem} Suppose $\mathcal{V}$ is a complex vector space and $T \in \mathcal{L}(\mathcal{V})$, then $T$ has an upper-triangular matrix with respect to some basis of $\mathcal{V}$.
\begin{proof} Note that result holds if $\text{dim}(\mathcal{V}) = 1$.\\
Suppose $\text{dim}(\mathcal{V}) > 1$ and the result holds for all complex vector spaces with dimension less than $\text{dim}(\mathcal{V})$. Let $\lambda$ be any eigenvalue of $T$, and let $\mathcal{U} = \text{range}(T - \lambda I)$. As $T - \lambda I$ is not surjective, $\text{dim}(\mathcal{U}) < \text{dim}(\mathcal{V})$, hence $\mathcal{U}$ is invariant under $T$. To show this, let $\mathbf{u} \in \mathcal{U}$, then
\begin{align}
T \mathbf{u} = (T - \lambda I)\mathbf{u} + \lambda \mathbf{u} \nonumber
\end{align}
$(T - \lambda I) \mathbf{u} \in \mathcal{U}$ by definition, and $\lambda \mathbf{u} \in \mathcal{U}$. Hence, equation shows $T \mathbf{u} \in \mathcal{U}$, and $\mathcal{U}$ is invariant under $T$.\\
Hence $T | _{\mathcal{U}}$ is an operator on $\mathcal{U}$. By induction, $\exists$ basis $(\mathbf{u}_1, \ldots \mathbf{u}_m)$ of $\mathcal{U}$ with respect to which $T|_{\mathcal{U}}$ has an upper-triangular matrix. Hence for each $j$,
\begin{align}
T \mathbf{u}_j = (T|_{\mathcal{U}})(\mathbf{u}_j) \in \text{span}(\mathbf{u}_1, \ldots, \mathbf{u}_j) \nonumber
\end{align}
Extend $(\mathbf{u}_1, \ldots, \mathbf{u}_m)$ to basis $(\mathbf{u}_1, \ldots, \mathbf{u}_m, \mathbf{v}_1, \ldots, \mathbf{v}_n)$ of $\mathcal{V}$. For each $k$,
\begin{align}
T \mathbf{v}_k = (T - \lambda I)\mathbf{v}_k + \lambda \mathbf{v}_k \nonumber
\end{align}
Definition of $\mathcal{U}$ shows that $(T - \lambda I)\mathbf{v}_k \in \mathcal{U} = \text{span}(\mathbf{u}_1, \ldots, \mathbf{u}_m)$. Hence,
\begin{align}
T \mathbf{v}_k \in \text{span}(\mathbf{u}_1, \ldots, \mathbf{u}_m, \mathbf{v}_1, \ldots, \mathbf{v}_k) \nonumber
\end{align}
Hence $T$ has an upper triangular matrix with respect to basis $(\mathbf{u}_1, \ldots, \mathbf{u}_m, \mathbf{v}_1, \ldots, \mathbf{v}_k)$.
\end{proof}
\end{theorem}

\begin{proposition} Suppose $T \in \mathcal{L}(\mathcal{V})$ has an upper-triangular matrix with respect to some basis of $\mathcal{V}$. Then $T$ is invertible if and only if all entries of the diagonal of upper-triangular matrix are non-zero.
\begin{proof} Let $(\mathbf{v}_1, \ldots, \mathbf{v}_n)$ be basis of $\mathcal{V}$ with respect to $T$ which has upper-triangular matrix
\begin{align}
\mathcal{M}(T,(\mathbf{v}_1, \ldots, \mathbf{v}_n)) =
\begin{bmatrix}
\lambda_1 & & * \\
& \lambda_2 & & \\
& & \ddots & \\
& & & \lambda_n
\end{bmatrix} \nonumber
\end{align} 
($\Rightarrow$) If $\lambda_1 = 0$, then $T \mathbf{v}_1 = \mathbf{0}$, hence $T$ is not invertible\\
Suppose $1 < k \leq n$ and $\lambda_k = 0$, then $T$ maps each of the vectors $\mathbf{v}_1, \ldots, \mathbf{v}_{k-1}$ into $\text{span}(\mathbf{v}_1, \ldots, \mathbf{v}_{k-1})$. As $\lambda_k = 0$, the matrix representation implies $T \mathbf{v}_k \in \text{span}(\mathbf{v}_1, \ldots, \mathbf{v}_{k-1})$. Hence, define a linear map
\begin{align}
S: \text{span}(\mathbf{v}_1, \ldots, \mathbf{v}_k) \rightarrow \text{span}(\mathbf{v}_1, \mathbf{v}_{k-1}) \nonumber
\end{align}
by $S \mathbf{v} = T \mathbf{v}$ for $\mathbf{v} \in \text{span}(\mathbf{v}_1, \ldots, \mathbf{v}_k)$. Then $S$ is $T$ restricted to $\text{span}(\mathbf{v}_1, \ldots, \mathbf{v}_k)$.\\
Note $\text{span}(\mathbf{v}_1, \ldots, \mathbf{v}_k)$ has dimension $k$ and $\text{span}(\mathbf{v}_1, \ldots, \mathbf{v}_{k-1})$ has dimension $k-1$ (as $(\mathbf{v}_1, \ldots, \mathbf{v}_n)$ is linearly independent). As $\text{span}(\mathbf{v}_1, \ldots, \mathbf{v}_k)$ has larger dimension than $\text{span}(\mathbf{v}_1, \ldots, \mathbf{v}_{k-1})$, no linear map from $\text{span}(\mathbf{v}_1, \ldots, \mathbf{v}_k)$ to $\text{span}(\mathbf{v}_1, \ldots, \mathbf{v}_{k-1})$ is injective, hence there exists nonzero vector $\mathbf{v} \in \text{span}(\mathbf{v}_1, \ldots, \mathbf{v}_{k})$ such that $S \mathbf{v} = \mathbf{0}$. Hence $T \mathbf{v} = \mathbf{0}$, and hence $T$ is not invertible.\\
($\Leftarrow$) Suppose $T$ is not invertible. Thus $T$ is not injective, and hence there exists a nonzero vector $\mathbf{v} \in \mathcal{V}$ such that $T \mathbf{v} = \mathbf{0}$. As $(\mathbf{v}_1, \ldots, \mathbf{v}_n)$ is a basis of $\mathcal{V}$, then $\mathbf{v} = a_1 \mathbf{v}_1 + \cdots + a_k \mathbf{v}_k$ where $a_1, \ldots, a_k \in \mathbb{F}$ and $a_k \neq 0$ (represent $\mathbf{v}$ as linear combination of $(\mathbf{v}_1, \ldots, \mathbf{v}_k)$ then choose $k$ to be largest index with nonzero coefficient).
\begin{align}
\mathbf{0} &= T \mathbf{v} = T(a_1 \mathbf{v}_1 + \cdots + a_k \mathbf{v}_k) \nonumber \\
&= (a_1 T \mathbf{v}_1 + \cdots + a_{k-1} T \mathbf{v}_{k-1}) + a_k T \mathbf{v}_k \nonumber
\end{align}
The last term in parentheses is in $\text{span}(\mathbf{v}_1, \ldots, \mathbf{v}_{k-1})$ (due to upper-triangular form).\\
Hence this shows $a_k T \mathbf{v}_k \in \text{span}(\mathbf{v}_1, \ldots, \mathbf{v}_{k-1})$. Multiply by $1/a_k$, which is allowed as $a_k \neq 0$, conclude that $T \mathbf{v}_k \in \text{span}(\mathbf{v}_1, \ldots, \mathbf{v}_{k-1})$. Thus when $T \mathbf{v}_k$ is written as linear combination of the basis $(\mathbf{v}_1, \ldots, \mathbf{v}_n)$, the coefficient of $\mathbf{v}_k$ will be $0$, hence $\lambda_k$ must be $0$, completing the proof.
\end{proof}
\end{proposition}

\begin{proposition} Suppose $T \in \mathcal{L}(\mathcal{V})$ has upper-triangular matrix with respect to some basis of $\mathcal{V}$, then eigenvalues of $T$ consists of precisely the entries of the diagonal of the upper-triangular matrix.
\begin{proof} Suppose $(\mathbf{v}_1, \ldots, \mathbf{v}_n)$ is a basis of $\mathcal{V}$ with respect to $T$ which $T$ has an upper-triangular matrix
\begin{align}
\mathcal{M}(T,(\mathbf{v}_1, \ldots, \mathbf{v}_n)) =
\begin{bmatrix}
\lambda_1 & & & * \\
& \lambda_2 & & \\
& & \ddots & \\
0 & & & \lambda_n
\end{bmatrix} \nonumber
\end{align} 
Let $\lambda \in \mathbb{F}$, then
\begin{align}
\mathcal{M}(T - \lambda I,(\mathbf{v}_1, \ldots, \mathbf{v}_n)) =
\begin{bmatrix}
\lambda_1 - \lambda & & & * \\
& \lambda_2 - \lambda & & \\
& & \ddots & \\
0 & & & \lambda_n - \lambda
\end{bmatrix} \nonumber
\end{align} 
Hence $T - \lambda I$ is not invertible if and only if $\lambda$ equals one of the $\lambda_j$.\\
Hence $\lambda$ is an eigenvalue of $T$ if and only if $\lambda$ equals one of the $\lambda_j$.
\end{proof}
\end{proposition}

\begin{definition} A \hlt{diagonal matrix} is a square matrix that is $0$ everywhere except possibly along diagonal.
\end{definition}

Operator $T \in \mathcal{L}(\mathcal{V})$ has diagonal matrix $\text{diag}(\lambda_1, \ldots, \lambda_n)$ w.r.t. basis $(\mathbf{v}_1, \ldots, \mathbf{v}_n)$ of $\mathcal{V}$ if and only if
\begin{align}
T \mathbf{v}_1 &= \lambda_1 \mathbf{v}_1 \nonumber \\
&\vdots \nonumber \\
T \mathbf{v}_n &= \lambda_n \mathbf{v}_n \nonumber
\end{align}
This follows from definition of matrix of operator with respect to a basis,\\

If an operator has a diagonal matrix w.r.t. some basis, the entires along diagonal are precisely eigenvalues.

\begin{proposition} If $T \in \mathcal{L}(\mathcal{V})$ has $\text{dim}(\mathcal{V})$ distinct eigenvalues, then $T$ has a diagonal matrix with respect to some basis of $\mathcal{V}$.
\begin{proof} Suppose $T \in \mathcal{L}(\mathcal{V})$ has $\text{dim}(\mathcal{V})$ distinct eigenvalues $\lambda_1, \ldots, \lambda_{\text{dim}(\mathcal{V})}$. For each $j$, let $\mathbf{v} \in \mathcal{V}$ be nonzero eigenvector corresponding to the eigenvalue $\lambda_j$. As nonzero eigenvectors corresponding to distinct eigenvalues are linearly independent, $(\mathbf{v}_1, \ldots, \mathbf{v}_{\text{dim}(\mathcal{V})})$ is linearly independent. A linearly independent list of $\text{dim}(\mathcal{V})$ vectors in $\mathcal{V}$ is a basis of $\mathcal{V}$, hence $(\mathbf{v}_1, \ldots, \mathbf{v}_{\text{dim}(\mathcal{V})})$ is a basis of $\mathcal{V}$. With respect to this basis, $T$ is a diagonal matrix.
\end{proof}
\end{proposition}

\begin{proposition} Suppose $T \in \mathcal{L}(\mathcal{V})$. Let $\lambda_1, \ldots, \lambda_m$ denote distinct eigenvalues of $T$. Then the following are equivalent:
\begin{enumerate}[label=\roman*.]
\setlength{\itemsep}{0pt}
\item $T$ has a diagonal matrix with respect to some basis of $\mathcal{V}$;
\item $\mathcal{V}$ has a basis consisting of eigenvectors of $T$;
\item $\exists$ one-dimensional subspaces $\mathcal{U}_1, \ldots, \mathcal{U}_n$ of $\mathcal{V}$, each invariant under $T$, such that
\begin{align}
\mathcal{V} = \mathcal{U}_1 \oplus \cdots \oplus \mathcal{U}_n \nonumber
\end{align}
\item $\mathcal{V} = \text{null}(T - \lambda_1 I) \oplus \cdots \oplus \text{null}(T - \lambda_m I)$
\item $\text{dim}(\mathcal{V}) = \text{dim}(\text{null}(T - \lambda_1 I)) + \cdots + \text{dim}(\text{null}(T - \lambda_m I))$
\end{enumerate}
\begin{proof} Already shown earlier that $(i)$ and $(ii)$ are equivalent. \\
$(ii) \Rightarrow (iii)$. Suppose $(ii)$ holds, thus $\mathcal{V}$ has a basis $(\mathbf{v}_1, \ldots, \mathbf{v}_n)$ consisting of eigenvectors of $T$. For each $j$, let $\mathcal{U}_j = \text{span}(\mathbf{v}_j)$, and hence each $\mathcal{U}_j$ is a one-dimensional subspace of $\mathcal{V}$ that is invariant under $T$ (as each $\mathbf{v}_j$ is an eigenvector of $T$). As $(\mathbf{v}_1, \ldots, \mathbf{v}_n)$ is a basis of $\mathcal{V}$, each vector in $\mathcal{V}$ can be written uniquely as a linearly combination of $(\mathbf{v}_1, \ldots, \mathbf{v}_n)$. Hence each vector in $\mathcal{V}$ can be written uniquely as a sum $\mathbf{u}_1 + \cdots + \mathbf{u}_n$, $\mathbf{u}_j \in \mathcal{U}_j$. Thus $\mathcal{V} = \mathcal{U}_1 \oplus \cdots \oplus \mathcal{U}_n$, and $(ii)$ implies $(iii)$.\\
$(iii) \Rightarrow (ii)$. Suppose $(iii)$ holds, hence there are one-dimensional subspaces $\mathcal{U}_1, \ldots, \mathcal{U}_n$ of $\mathcal{V}$, each invariant under $T$, such that $\mathcal{V} = \mathcal{U}_1 \oplus \cdots \oplus \mathcal{U}_n$. For each $j$, let $\mathbf{v}_j$ be nonzero vector in $\mathcal{U}_j$. Then each $\mathbf{v}_j$ is an eigenvector of $T$. As each vector in $\mathcal{V}$ can be written uniquely as $\mathbf{u}_1 + \cdots + \mathbf{u}_n, \mathbf{u}_j \in \mathcal{U}_j$, $(\mathbf{v}_1, \ldots, \mathbf{v}_n)$ is a basis of $\mathcal{V}$.\\
Hence this shows so far $(i) \Leftrightarrow (ii) \Leftrightarrow (iii)$.\\
$(ii) \Rightarrow (iv)$. Suppose $(ii)$ holds, then $\mathcal{V}$ has a basis consisting of eigenvectors of $T$, and hence every vector in $\mathcal{V}$ is a linear combination of eigenvectors of $T$. Thus 
\begin{align}
\mathcal{V} = \text{null}(T - \lambda_1 I) + \cdots + \text{null}(T - \lambda_m I) \nonumber
\end{align}
Suppose $\mathbf{0} = \mathbf{u}_1 + \cdots + \mathbf{u}_m, \mathbf{u}_j \in \text{null}(T - \lambda_j I)$. As nonzero eigenvectors corresponding to distinct eigenvalues are linearly independent, this implies each $\mathbf{u}_j = \mathbf{0}$. This implies the sum is a direct sum.\\
$(iv) \Rightarrow (v)$. Suppose $(iv)$ holds. By definition, each $\text{null}(T - \lambda_j I)$ is a subspace of $\mathcal{V}$. By definition of direct sum, each $\mathbf{V} \in \mathcal{V}$ can be uniquely expressed as $\mathbf{v} = \mathbf{u}_1 + \cdots + \mathbf{u}_m, \mathbf{u}_j \in \text{null}(T - \lambda_j I)$.\\
The only way for each $\mathbf{u}_1 + \cdots + \mathbf{u}_m = \mathbf{0}$ with each $\mathbf{u}_j \in \text{null}(T - \lambda_j I)$ is if each $\mathbf{u}_j = \mathbf{0}$. Hence these properties together imply $(iv) \Rightarrow (v)$.\\
$(v) \Rightarrow (ii)$. Suppose $(v)$ holds. Choose a basis $\mathbf{v}_j \in \text{null}(T - \lambda_j I)$, and form a list $(\mathbf{v}_1, \ldots, \mathbf{v}_n)$ of eigenvectors of $T$, where $n = \text{dim}(\mathcal{V})$. Suppose $a_1 \mathbf{v}_1 + \cdots + a_n \mathbf{v}_n = \mathbf{0}$, where $a_1, \ldots, a_n \in \mathbb{F}$. For each $j = 1, \ldots, m$, let $\mathbf{u}_j$ be sum of all terms $a_k \mathbf{v}_k$ such that $\mathbf{v}_k \in \text{null}(T - \lambda_j I)$. Thus each $\mathbf{u}_j$ is an eigenvector of $T$ with eigenvalue $\lambda_j$, and $\mathbf{u}_1 + \cdots + \mathbf{u}_m = \mathbf{0}$. As nonzero eigenvectors corresponding to distinct eigenvalues are linearly independent, this implies each $\mathbf{u}_j = \mathbf{0}$. As each $\mathbf{u}_j$ is sum of $a_k \mathbf{v}_k$, where $\mathbf{v}_k$ were chosen to be a basis of $\text{null}(T - \lambda_j I)$, this implies all $a_k = 0$. Hence $(\mathbf{v}_1, \ldots, \mathbf{v}_n)$ is linearly independent and hence is a basis of $\mathcal{V}$.
\end{proof}
\end{proposition}

Next theorem shows that an invariant subspace of dimension $1$ or $2$ always exists.

\begin{theorem} \label{thm:invariantsubspaceexistence} Every operator on a finite-dimensional, nonzero, real vector space has an invariant subspace of dimension $1$ or $2$.
\begin{proof} Let $\mathcal{V}$ be a real vector space, $\text{dim}(\mathcal{V}) = n > 0$, and $T \in \mathcal{L}(\mathcal{V})$. Choose $\mathbf{v} \in \mathcal{V}$ with $\mathbf{v} \neq \mathbf{0}$, then
\begin{align}
(\mathbf{v}, T\mathbf{v}, T^2 \mathbf{v}, \ldots, T^n \mathbf{v}) \nonumber
\end{align}
cannot be linearly independent as $\mathcal{V}$ has dimension $n$, and there are $n+1$ vectors.\\
Hence $\exists \ a_0, \ldots a_n \in \R$ not all $0$, such that
\begin{align}
\mathbf{0} = a_0 \mathbf{v} + a_1 T \mathbf{v} + \cdots + a_n T^n \mathbf{v} \nonumber
\end{align}
Let $a$ be coefficients of a polynomial, then the factored form is
\begin{align}
a_0 + a_1 x + \cdots + a_n x^n = c(x - \lambda_1) \cdots (x - \lambda_m)(x^2 \alpha_1 x + \beta_1) \cdots (x^2 + \alpha_M x + \beta_M) \nonumber
\end{align}
where $c \neq 0 \in \R$, $\lambda_j, \alpha_j, \beta_j \in \R$, $m + M \geq 1$, and equation holds for all $x \in \R$. Then
\begin{align}
\mathbf{0} &= a_0 \mathbf{v} + a_1 T \mathbf{v} + \cdots + a_n T^n \mathbf{v} \nonumber \\
&= (a_0 I + a_1 T + \cdots + a_n T^n)\mathbf{v} \nonumber \\
&= c(T - \lambda_1 I)\cdots (T - \lambda_m I)(T^2 + \alpha_1 T + \beta_1 I) \cdots (T^2 + \alpha_M T + \beta_M I)\mathbf{v} \nonumber
\end{align}
This means $T - \lambda_j I$ is not injective for at least one $j$, or $(T^2 + \alpha_j T + \beta_j I)$ is not injective for at least one $j$. If $T - \lambda_j I$ is not injective for at least one $j$, then $T$ has an eigenvalue, and hence a one-dimensional subspace.\\
If $(T^2 + \alpha_j T + \beta_j T)$ is not injective for some $j$, then there exists nonzero vector $\mathbf{u} \in \mathcal{V}$ such that
\begin{align}
T^2 \mathbf{u} + \alpha_j T \mathbf{u} + \beta_j \mathbf{u} = \mathbf{0} \nonumber
\end{align}
Consider typical element of $\text{span}(\mathbf{u}, T\mathbf{u})$ of form $a \mathbf{u} + b T \mathbf{u}$, where $a, b \in \R$, then
\begin{align}
T(a \mathbf{u} + bT \mathbf{u}) &= a T \mathbf{u} + b T^2 \mathbf{u} \nonumber \\
&= a T \mathbf{u} - b \alpha_j T \mathbf{u} - b \beta_j \mathbf{u} \nonumber
\end{align}
last equality comes from solving the equality for $T^2 \mathbf{u}$. This shows $T(a \mathbf{u} + b T \mathbf{u}) \in \text{span}(\mathbf{u}, T \mathbf{u})$. hence $\text{span}(\mathbf{u}, T \mathbf{u})$ is invariant under $T$.
\end{proof}
\end{theorem}

\begin{definition} Let $\mathcal{U}$, $\mathcal{W}$ be subspaces of $\mathcal{V}$, where $\mathcal{V} = \mathcal{U} \oplus \mathcal{W}$.\\
Let $\mathbf{u} \in \mathcal{U}, \mathbf{w} \in \mathcal{W}, \mathbf{v} \in \mathcal{V}$, and $\mathbf{v} = \mathbf{u} + \mathbf{w}$. Define \hlt{projection} $P_{\mathcal{U}, \mathcal{W}} \in \mathcal{L}(\mathcal{V})$ by $P_{\mathcal{U}, \mathcal{W}} \mathbf{v} = \mathbf{u}$. Note that
\begin{align}
P_{\mathcal{W}, \mathcal{U}} \mathbf{v} &= \mathbf{w}, \ \ \ \mathbf{v} = P_{\mathcal{U}, \mathcal{W}} \mathbf{v} + P_{\mathcal{W}, \mathcal{U}} \mathbf{v} \nonumber \\
\text{range} (P_{\mathcal{U}, \mathcal{W}}) &= \mathcal{U}, \ \ \ \text{null}(P_{\mathcal{U}, \mathcal{W}}) = \mathcal{W} \nonumber
\end{align}
\end{definition}

\begin{theorem} Every operator on an odd-dimensional real vector space has an eigenvalue.
\begin{proof} Suppose $\mathcal{V}$ is a real vector space with odd dimension.\\
Note that result holds for $\text{dim}(\mathcal{V}) = 1$.\\
Suppose $\text{dim}(\mathcal{V})$ is an odd number greater than $1$. Suppose $T \in \mathcal{L}(\mathcal{V})$, to prove $T$ has an eigenvalue. By Theorem \ref{thm:invariantsubspaceexistence}, there is a two-dimensional subspace $\mathcal{U}$ of $\mathcal{V}$ that is invariant under $T$. Let $\mathcal{W}$ be any subspace of $\mathcal{V}$ such that $\mathcal{V} = \mathcal{U} \oplus \mathcal{W}$, then Theorem \ref{thm:invariantsubspaceexistence} guarantees that such a $\mathcal{W}$ exists.\\
As $\mathcal{W}$ has dimension $2$ less than $\text{dim}(\mathcal{V})$, apply induction hypothesis to $T|_{\mathcal{W}}$. However, $\mathcal{W}$ might not be invariant under $T$, which means $T_{\mathcal{W}}$ might not be an operator on $\mathcal{W}$.\\
Define $S \in \mathcal{L}(\mathcal{W})$ by $S\mathbf{w} = P_{\mathcal{W}, \mathcal{U}} (T \mathbf{w})$ for $\mathbf{w} \in \mathcal{W}$. By induction, $S$ has an eigenvalue $\lambda$.\\
Let $\mathbf{w} \in \mathcal{W}$ be nonzero eigenvector for $S$ corresponding to the eigenvalue $\lambda$, thus $(S - \lambda I)\mathbf{w} = 0$. To look for an eigenvector of $T$ in $\mathcal{U} + \text{span}(\mathbf{w})$. Consider typical vector $\mathbf{u} + a \mathbf{w}$ in $\mathcal{U} = \text{span}(\mathbf{w})$, where $\mathbf{u} \in \mathcal{U}$ and $a \in \R$.
\begin{align}
(T - \lambda I)(\mathbf{u} + a \mathbf{w}) &= T \mathbf{u} - \lambda \mathbf{u} + a (T \mathbf{w} - \lambda \mathbf{w}) \nonumber \\
&= T \mathbf{u} - \lambda \mathbf{u} + a (P_{\mathcal{U}, \mathcal{W}}(T \mathbf{w}) + P_{\mathcal{W}, \mathcal{U}}(T \mathbf{w}) - \lambda \mathbf{w}) \nonumber \\
&= T \mathbf{u} - \lambda \mathbf{u} + a (P_{\mathcal{U}, \mathcal{W}}(T \mathbf{w}) + S \mathbf{w} - \lambda \mathbf{w}) \nonumber \\
&= T \mathbf{u} - \lambda \mathbf{u} + a P_{\mathcal{U}, \mathcal{W}}(T \mathbf{w}) \nonumber
\end{align}
Note $T \mathbf{u} \in \mathcal{U}$ (as $\mathcal{U}$ is invariant under $T$), $\lambda \mathbf{u} \in \mathcal{U}$ (as $\mathbf{u} \in \mathcal{U}$), and $a P_{\mathcal{U}, \mathcal{W}}(T \mathbf{w}) \in \mathcal{U}$ (by definition). Hence $T - \lambda I$ maps $\mathcal{U}+ \text{span}(\mathbf{w})$ into  $\mathcal{U}$. As $\mathcal{U}+ \text{span}(\mathbf{w})$ has larger dimension than $\mathcal{U}$ this means $(T - \lambda I)|_{\mathcal{U}+ \text{span}(\mathbf{w})}$ is not injective, and there exists a nonzero vector $\mathbf{v} \in \mathcal{U}+ \text{span}(\mathbf{w})$ such that $(T - \lambda I) \mathbf{v} = \mathbf{0}$. Hence $T$ has an eigenvalue as desired.
\end{proof}
\end{theorem}

\subsubsection{Inner-Product Spaces}

\begin{definition} The \hlt{norm} is defined as $\norm{\mathbf{z}} = \sqrt{\abs{z_1}^2 + \cdots + \abs{z_n}^2} = \sqrt{\innerprod{\mathbf{z}, \mathbf{z}}}$ for $\mathbf{z} = (z_1, \ldots, z_n) \in \C^n$.\\
The norm should be think of as the \hlt{inner product} of $\mathbf{z}$ with itself, $\norm{z}^2 = z_1 \overline{z_1} + \cdots + z_n \overline{z_n}$.\\
If $\mathbf{p} \in \mathcal{P}_m(\mathbb{F})$ (polynomial space0, then $\norm{\mathbf{p}} = \sqrt{\int_0^1 \abs{\mathbf{p}(x)}^2 \ dx}$
\end{definition}

\begin{remark} An inner product on $\mathcal{V}$ that is a function that takes each ordered pair $(\mathbf{u}, \mathbf{v})$ of elements in $\mathcal{V}$ to a number $\innerprod{\mathbf{u}, \mathbf{v}} \in \mathbb{F}$ has following properties:
\begin{enumerate}[label=\roman*.]
\setlength{\itemsep}{0pt}
\item Positivity: $\innerprod{\mathbf{v}, \mathbf{v}} \geq 0 \ \forall \ \mathbf{v} \in \mathcal{V}$
\item Definiteness: $\innerprod{\mathbf{v}, \mathbf{v}} = 0$ if and only if $\mathbf{v} = \mathbf{0}$
\item Additivity in First Slot: $\innerprod{\mathbf{u} + \mathbf{v}, \mathbf{w}} = \innerprod{\mathbf{u}, \mathbf{w}} + \innerprod{\mathbf{v}, \mathbf{w}}$ for all $\mathbf{u}, \mathbf{v}, \mathbf{w} \in \mathcal{V}$
\item Homogeneity in First Slot: $\innerprod{a\mathbf{v}, \mathbf{w}} = a \innerprod{\mathbf{v}, \mathbf{w}}$ for all $a \in \mathbb{F}$ and $\mathbf{v}, \mathbf{w} \in \mathcal{V}$
\item Conjugate Symmetry: $\innerprod{\mathbf{v}, \mathbf{w}} = \overline{\innerprod{\mathbf{w}, \mathbf{v}}}$ for all $\mathbf{v}, \mathbf{w} \in \mathcal{V}$
\end{enumerate}
\end{remark}

\begin{definition} An \hlt{inner-product space} is a vector space $\mathcal{V}$ with an inner product on $\mathcal{V}$.
\end{definition}

\begin{definition} Two vectors $\mathbf{u}$, $\mathbf{v}$ are \hlt{orthogonal} if $\innerprod{\mathbf{u}, \mathbf{v}} = 0$.
\end{definition}

\begin{theorem} \label{thm:pythagorean} \hlt{Pythagorean Theorem}\\
If $\mathbf{u}, \mathbf{v}$ are orthogonal vectors in $\mathcal{V}$, then $\norm{\mathbf{u} + \mathbf{v}}^2 = \norm{\mathbf{u}}^2 + \norm{\mathbf{v}}^2$.
\begin{proof} Suppose $\mathbf{u}, \mathbf{v}$ are orthogonal vectors in $\mathcal{V}$, then
\begin{align}
\norm{\mathbf{u} + \mathbf{v}}^2 &= \innerprod{\mathbf{u} + \mathbf{v}, \mathbf{u} + \mathbf{v}} \nonumber \\
&= \norm{\mathbf{u}}^2 + \norm{\mathbf{u}}^2 + \innerprod{\mathbf{u}, \mathbf{v}} + \innerprod{\mathbf{v}, \mathbf{u}} \nonumber \\
&= \norm{\mathbf{u}}^2 + \norm{\mathbf{v}}^2 \nonumber
\end{align}
\end{proof}
\end{theorem}

\begin{figure}[h]
  \centering
  \begin{tikzpicture}[scale=2,>=stealth,
  point/.style={circle,fill,inner sep=0.9pt},
  every node/.style={font=\small}
]
  \coordinate (O) at (0,0);
  \coordinate (D) at (1,0);
  \coordinate (U) at (1.6,1.0);
  \coordinate (P) at ($(O)!(U)!(D)$);
  \coordinate (V) at (1,0);

  \draw[thin] ($(O)!-0.2!(D)$) -- ($(O)!1.8!(D)$);
  \draw[->] (O) -- (U) node[above] {$\mathbf{u}$};
  \draw (U) -- (P) node[midway,right] {$\mathbf{w}$};
  \node[point,label=below left:{$\mathbf{0}$}] at (O) {};
  \node[point,label=below:{$\mathbf{v}$}] at (V) {};
  \node[point,label=below:{$\lambda \mathbf{v}$}] at (P) {};
\end{tikzpicture}
  \caption{Orthogonal decomposition of $\mathbf{u}$ as a scalar multiple of $\mathbf{v}$ plus vector $\mathbf{w}$ orthogonal to $\mathbf{v}$}
  \label{fig:orthodecomp}
\end{figure}

\begin{definition} To compute the \hlt{orthogonal decomposition} of vector $\mathbf{u}$, let $a \in \mathbb{F}$, then
\begin{align}
\mathbf{u} = a \mathbf{v} + (\mathbf{u} - a \mathbf{v}) \nonumber
\end{align}
To choose $a$ so that $\mathbf{v}$ is orthogonal to $(\mathbf{u} - a \mathbf{v})$, which is
\begin{align}
0 = \innerprod{\mathbf{u} - a \mathbf{v}, \mathbf{v}} = \innerprod{\mathbf{u}, \mathbf{v}} - a \norm{\mathbf{v}}^2 \nonumber
\end{align}
Need to choose $a$ to be $\innerprod{\mathbf{u}, \mathbf{v}}/\norm{\mathbf{v}}^2$, where $\mathbf{v} \neq \mathbf{0}$, then
\begin{align}
\mathbf{u} = \frac{\innerprod{\mathbf{u}, \mathbf{v}}}{\norm{\mathbf{v}}^2} \mathbf{v} + \left( \mathbf{u} - \frac{\innerprod{\mathbf{u}, \mathbf{v}}}{\norm{\mathbf{u}}^2}\mathbf{v} \right) \nonumber
\end{align}
\end{definition}

\begin{theorem} \label{thm:cauchyschwarzineq} \hlt{Cauchy-Schwarz Inequality}\\
If $\mathbf{u}, \mathbf{v} \in \mathcal{V}$, then $\abs{\innerprod{\mathbf{u}, \mathbf{v}}} \leq \norm{\mathbf{u}} \norm{\mathbf{u}}$.\\
This inequality is an equality if and only if one of $\mathbf{u}, \mathbf{v}$ is a scalar multiple of the other.
\begin{proof} Let $\mathbf{u}, \mathbf{v} \in \mathcal{V}$. If $\mathbf{v} = \mathbf{0}$, then inequality holds.\\
Assume $\mathbf{u} \neq \mathbf{0}$. Consider orthogonal decomposition
\begin{align}
\mathbf{u} = \frac{\innerprod{\mathbf{u}, \mathbf{v}}}{\norm{\mathbf{v}}^2} \mathbf{v} + \mathbf{w} \nonumber
\end{align}
where $\mathbf{w}$ is orthogonal to $\mathbf{v}$. By Pythagorean Theorem \ref{thm:pythagorean},
\begin{align}
\norm{\mathbf{u}}^2 &= \norm{\frac{\innerprod{\mathbf{u}, \mathbf{v}}}{\norm{\mathbf{v}}^2} \mathbf{v}}^2 + \norm{\mathbf{w}}^2 \nonumber \\
&= \frac{\abs{\innerprod{\mathbf{u}, \mathbf{v}}}^2}{\norm{\mathbf{v}}^2} + \norm{\mathbf{w}}^2 \nonumber \\
&\geq \frac{\abs{\innerprod{\mathbf{u}, \mathbf{v}}}^2}{\norm{\mathbf{v}}^2} \nonumber
\end{align}
Multiplying both sides by $\norm{\mathbf{v}}^2$ then taking square roots give the desired inequality.\\
Note that this is an equality if and only if $\mathbf{w} = \mathbf{0}$, but $\mathbf{w} = \mathbf{0}$ if and only if $\mathbf{u}$ is a multiple of $\mathbf{v}$.
\end{proof}
\end{theorem}

The next result is on triangle inequality, which states the length of any side of a triangle is less than the sum of length of the other two sides.

\begin{theorem} \label{thm:triangleineq} \hlt{Triangle Inequality}\\
If $\mathbf{u}, \mathbf{v} \in \mathcal{V}$, then $\norm{\mathbf{u} + \mathbf{v}} \leq \norm{\mathbf{u}} + \norm{\mathbf{v}}$.\\
This inequality is an equality if and only if one of $\mathbf{u}, \mathbf{v}$ is a nonnegative multiple of the other.
\begin{proof} Let $\mathbf{u}, \mathbf{v} \in \mathcal{V}$, then
\begin{align}
\norm{\mathbf{u} + \mathbf{v}}^2 &= \innerprod{\mathbf{u} + \mathbf{v}, \mathbf{u} + \mathbf{v}} \nonumber \\
&= \innerprod{\mathbf{u}, \mathbf{u}} + \innerprod{\mathbf{v}, \mathbf{v}} + \innerprod{\mathbf{u}, \mathbf{v}} + \innerprod{\mathbf{v}, \mathbf{u}} \nonumber \\
&= \innerprod{\mathbf{u}, \mathbf{u}} + \innerprod{\mathbf{v}, \mathbf{v}} + \innerprod{\mathbf{u}, \mathbf{v}} + \overline{\innerprod{\mathbf{u}, \mathbf{v}}} \nonumber \\
&= \norm{\mathbf{u}}^2 + \norm{\mathbf{u}}^2 + 2 \text{Re}\innerprod{\mathbf{u}, \mathbf{v}} \nonumber \\
&\leq \norm{\mathbf{u}}^2 + \norm{\mathbf{v}}^2 + 2 \abs{\innerprod{\mathbf{u}, \mathbf{v}}} \nonumber \\
&\leq \norm{\mathbf{u}}^2 + \norm{\mathbf{v}}^2 + 2 \norm{\mathbf{u}} \norm{\mathbf{v}} \ \ \ \ \ \text{by Cauchy-Schwarz Inequality Theorem \ref{thm:cauchyschwarzineq}} \nonumber \\
&= (\norm{\mathbf{u}} + \norm{\mathbf{v}})^2 \nonumber
\end{align}
Taking square roots on both sides gives the triangle inequality.\\
Note this is an equality if and only if $\innerprod{\mathbf{u}, \mathbf{v}} = \norm{\mathbf{u}} \norm{\mathbf{v}}$.
\end{proof}
\end{theorem}

The next result is on parallelogram inequality, which states that in any parallelogram, the sum of squares of length of diagonals equals the sum of lengths of the four sides.

\begin{theorem} \label{thm:parallelogramineq} \hlt{Parallelogram Inequality}\\
If $\mathbf{u}, \mathbf{v} \in \mathcal{V}$, then $\norm{\mathbf{u} + \mathbf{v}}^2 + \norm{\mathbf{u} - \mathbf{v}}^2 = 2 (\norm{\mathbf{u}}^2 + \norm{\mathbf{v}}^2)$.
\begin{proof} Let $\mathbf{u}, \mathbf{v} \in \mathcal{V}$, then
\begin{align}
\norm{\mathbf{u} + \mathbf{v}}^2 + \norm{\mathbf{u} - \mathbf{v}}^2 &= \innerprod{\mathbf{u} + \mathbf{v}, \mathbf{u} + \mathbf{v}} + \innerprod{\mathbf{u} - \mathbf{v}, \mathbf{u} - \mathbf{v}} \nonumber \\
&= 2\norm{\mathbf{u}}^2 + 2\norm{\mathbf{v}}^2 + \innerprod{\mathbf{u}, \mathbf{v}} + \innerprod{\mathbf{v}, \mathbf{u}} - \innerprod{\mathbf{u}, \mathbf{v}} - \innerprod{\mathbf{v}, \mathbf{u}} \nonumber \\
&= 2 (\norm{\mathbf{u}}^2 + \norm{\mathbf{v}}^2) \nonumber
\end{align}
\end{proof}
\end{theorem} 

\begin{definition} A list of vectors $(\mathbf{e}_1, \ldots, \mathbf{e}_m)$ in $\mathcal{V}$ is \hlt{orthonormal} if 
\begin{align}
\innerprod{\mathbf{e}_j, \mathbf{e}_k} =
\begin{cases}
\innerprod{\mathbf{e}_j, \mathbf{e}_k} = 0 &\text{if } j \neq k \\
\innerprod{\mathbf{e}_j, \mathbf{e}_k} = 1 &\text{if } j = k
\end{cases} \nonumber
\end{align}
\end{definition}

\begin{proposition} \label{prop:sumoforthonormal} If $(\mathbf{e}_1, \ldots, \mathbf{e}_m)$ is an orthonormal list of vectors in $\mathcal{V}$, then
\begin{align}
\norm{a_1 \mathbf{e}_1 + \cdots + a_m \mathbf{e}_m}^2 = \abs{a_1}^2 + \cdots + \abs{a_m}^2 \nonumber
\end{align}
for all $a_1, \ldots, a_m \in \mathbb{F}$.
\begin{proof} As each $\mathbf{e}_j$ has norm $1$, by Pythagorean Theorem \ref{thm:pythagorean}, this holds.
\end{proof}
\end{proposition}

\begin{corollary} \label{cor:orthonormallistisindep} Every orthonormal list of vectors is linearly independent.
\begin{proof} Suppose $(\mathbf{e}_1, \ldots, \mathbf{e}_m)$ is an orthonormal list of vectors in $\mathcal{V}$ and $a_1, \ldots, a_m \in \mathbb{F}$ are such that
\begin{align}
a_1 \mathbf{e}_1 + \cdots + a_m \mathbf{e}_m = \mathbf{0} \nonumber
\end{align}
then $\abs{a_1}^2 + \cdots + \abs{a_m}^2 = 0$ by Proposition \ref{prop:sumoforthonormal}, which means all $a_j$ are $0$ as desired.
\end{proof}
\end{corollary}

\begin{definition} An \hlt{orthonormal basis} of $\mathcal{V}$ is an orthonormal list of vectors in $\mathcal{V}$ that is also a basis.
\end{definition}

\begin{theorem} Suppose $(\mathbf{e}_1, \ldots, \mathbf{e}_n)$ is an orthonormal basis of $\mathcal{V}$, then
\begin{align}
\mathbf{v} &= \innerprod{\mathbf{v}, \mathbf{e}_1}\mathbf{e}_1 + \cdots + \innerprod{\mathbf{v}, \mathbf{e}_n}\mathbf{e}_n \nonumber \\
\norm{\mathbf{v}}^2 &= \abs{\innerprod{\mathbf{v}, \mathbf{e}_1}}^2 + \cdots + \abs{\innerprod{\mathbf{v}, \mathbf{e}_n}}^2 \nonumber
\end{align}
for every $\mathbf{e} \in \mathcal{V}$.
\begin{proof} Let $\mathbf{v} \in \mathcal{V}$. As $(\mathbf{e}_1, \ldots, \mathbf{e}_n)$ is a basis of $\mathcal{V}$, there exists scalars $a_1, \ldots, a_n$ such that
\begin{align}
\mathbf{v} = a_1 \mathbf{e}_1 + \cdots + a_n \mathbf{e}_n \nonumber
\end{align}
Take inner product of both sides of equation with $\mathbf{e}_j$ to get $\innerprod{\mathbf{v}, \mathbf{e}_j} = a_j$, hence proving.
\end{proof}
\end{theorem}

\begin{theorem} \label{thm:gramschmidtprocedure} If $(\mathbf{v}_1, \ldots, \mathbf{v}_m)$ is a linearly independent list of vectors in $\mathcal{V}$, then there exists an orthonormal list $(\mathbf{e}_1, \ldots, \mathbf{e}_m)$ of vectors in $\mathcal{V}$ such that
\begin{align}
\text{span}(\mathbf{v}_1, \ldots, \mathbf{v}_j) = \text{span}(\mathbf{e}_1, \ldots, \mathbf{e}_j) \nonumber
\end{align}
for $j = 1, \ldots, m$.
\begin{proof} Suppose $(\mathbf{v}_1, \ldots, \mathbf{v}_m)$ is linearly independent list of vectors in $\mathcal{V}$. To construct $\mathbf{e}$, let $\mathbf{e} = \mathbf{v}_1/\norm{\mathbf{v}_1}$, which satisfies condition for $j=1$.\\
Suppose $j > 1$, and orthonormal list $(\mathbf{e}_1, \ldots, \mathbf{e}_{j-1})$ has been chosen so that
\begin{align}
\text{span}(\mathbf{v}_1, \ldots, \mathbf{v}_{j-1}) = \text{span}(\mathbf{e}_1, \ldots, \mathbf{e}_{j-1}) \nonumber
\end{align}
Let
\begin{align}
\mathbf{e}_j = \frac{\mathbf{v}_j - \innerprod{\mathbf{v}_j, \mathbf{e}_1}\mathbf{e}_1 - \cdots - \innerprod{\mathbf{v}_j, \mathbf{e}_{j-1}}\mathbf{e}_{j-1}}{\norm{\mathbf{v}_j - \innerprod{\mathbf{v}_j, \mathbf{e}_1}\mathbf{e}_1 - \cdots - \innerprod{\mathbf{v}_j, \mathbf{e}_{j-1}}\mathbf{e}_{j-1}}} \nonumber
\end{align}
Note that $\mathbf{v}_j \notin \text{span}(\mathbf{v}_1, \ldots, \mathbf{v}_{j-1})$ (as $(\mathbf{v}_1, \ldots, \mathbf{v}_m)$ is linearly independent), hence $\mathbf{v}_j \notin \text{span}(\mathbf{e}_1, \ldots, \mathbf{e}_{j-1})$. Hence above equation is not dividing by $\mathbf{0}$, and $\mathbf{e}_j$ is well defined. Note $\norm{\mathbf{e}_j} = 1$. Let $1 \leq k < j$, then
\begin{align}
\innerprod{\mathbf{e}_j, \mathbf{e}_k} &= \innerprod{\frac{\mathbf{v}_j - \innerprod{\mathbf{v}_j, \mathbf{e}_1}\mathbf{e}_1 - \cdots - \innerprod{\mathbf{v}_j, \mathbf{e}_{j-1}}\mathbf{e}_{j-1}}{\norm{\mathbf{v}_j - \innerprod{\mathbf{v}_j, \mathbf{e}_1}\mathbf{e}_1 - \cdots - \innerprod{\mathbf{v}_j, \mathbf{e}_{j-1}}\mathbf{e}_{j-1}}}, \mathbf{e}_{k}} \nonumber \\
&= \frac{\innerprod{\mathbf{v}_j, \mathbf{e}_k} - \innerprod{\mathbf{v}_j, \mathbf{e}_k}}{\norm{\mathbf{v}_j - \innerprod{\mathbf{v}_j, \mathbf{e}_1}\mathbf{e}_1 - \cdots - \innerprod{\mathbf{v}_j, \mathbf{e}_{j-1}}\mathbf{e}_{j-1}}} = 0 \nonumber
\end{align}
Hence $(\mathbf{e}_1, \ldots, \mathbf{e}_j)$ is an orthonormal list.\\
Note that $\mathbf{v}_j \in \text{span}(\mathbf{e}_1, \ldots, \mathbf{e}_j)$, and hence
\begin{align}
\text{span}(\mathbf{v}_1, \ldots, \mathbf{v}_j) \subset \text{span}(\mathbf{e}_1, \ldots, \mathbf{e}_j) \nonumber
\end{align}
Both lists are linearly independent, hence both spaces have dimension $j$, thus must be equal.
\end{proof}
\end{theorem}

\begin{corollary} Every finite-dimensional inner-product space has an orthonormal basis.
\begin{proof} Choose a basis of $\mathcal{V}$. Apply Gram-Schmidt \ref{thm:gramschmidtprocedure} to it to produce an orthonormal list. This orthonormal list is linearly independent by Corollary \ref{cor:orthonormallistisindep} and its span equals $\mathcal{V}$, hence it is an orthonormal basis of $\mathcal{V}$.
\end{proof}
\end{corollary}










